\documentclass{article}%
\usepackage[T1]{fontenc}%
\usepackage[utf8]{inputenc}%
\usepackage{lmodern}%
\usepackage{textcomp}%
\usepackage{lastpage}%
\usepackage{geometry}%
\geometry{margin=1in,headheight=14pt,headsep=25pt}%
\usepackage{hyperref}%
\usepackage{enumitem}%
\usepackage{fancyhdr}%
\usepackage{xcolor}%
%

            \hypersetup{
                colorlinks=true,
                linkcolor=blue,
                filecolor=magenta,
                urlcolor=blue,
            }
            
            \pagestyle{fancy}
            \fancyhf{}
            \rhead{Generated on \today}
            \lhead{Course Summary}
            \cfoot{\thepage}
        %
\title{Course Summary}%
\author{Generated by ScribeLec}%
\date{\today}%
%
\begin{document}%
\normalsize%
\maketitle%
\tableofcontents\newpage%
\section{Key Terms}%
\label{sec:KeyTerms}%
\begin{itemize}[[leftmargin=*]]%
\item \href{https://www.math.purdue.edu/~yipn/421/2024-08-27-ExSimplex.pdf#page=1}{\textbf{Simplex Method}}: An iterative algorithm used to solve linear programming problems by moving from one vertex of the feasible region to another until an optimal solution is found.<L[2024-08-27-ExSimplex] 1><L[2024-08-27-ExSimplex] 6>%
\item \href{https://www.math.purdue.edu/~yipn/421/2024-08-27-ExSimplex.pdf#page=1}{\textbf{Slack Variable}}: A non-negative variable added to a less-than-or-equal-to constraint to transform it into an equality constraint.<L[2024-08-27-ExSimplex] 1>%
\item \href{https://www.math.purdue.edu/~yipn/421/2024-08-27-ExSimplex.pdf#page=2}{\textbf{Feasible Region}}: The set of all points that satisfy all the constraints of a linear programming problem.<L[2024-08-27-ExSimplex] 2><L[2024-08-27-ExSimplex] 4><L[2024-08-27-ExSimplex] 5><L[2024-08-27-ExSimplex] 16><L[2024-08-27-ExSimplex] 30>%
\item \href{https://www.math.purdue.edu/~yipn/421/2024-08-27-ExSimplex.pdf#page=4}{\textbf{Vertex}}: A point where two or more constraint boundaries intersect in the feasible region of a linear programming problem. Vertices represent potential optimal solutions.<L[2024-08-27-ExSimplex] 4><L[2024-08-27-ExSimplex] 5><L[2024-08-27-ExSimplex] 16><L[2024-08-27-ExSimplex] 30>%
\item \href{https://www.math.purdue.edu/~yipn/421/2024-08-27-ExSimplex.pdf#page=9}{\textbf{Basic Variable}}: In the Simplex method, a variable that is expressed in terms of non-basic variables in a dictionary.<L[2024-08-27-ExSimplex] 9><L[2024-08-27-ExSimplex] 10><L[2024-08-27-ExSimplex] 26><L[2024-08-27-ExSimplex] 27><L[2024-08-27-ExSimplex] 28>%
\item \href{https://www.math.purdue.edu/~yipn/421/2024-08-27-ExSimplex.pdf#page=6}{\textbf{Non{-}basic Variable}}: In the Simplex method, a variable set to zero in a given iteration of the algorithm.<L[2024-08-27-ExSimplex] 6><L[2024-08-27-ExSimplex] 7><L[2024-08-27-ExSimplex] 26><L[2024-08-27-ExSimplex] 27><L[2024-08-27-ExSimplex] 28>%
\item \href{https://www.math.purdue.edu/~yipn/421/2024-08-27-ExSimplex.pdf#page=9}{\textbf{Pivot Operation}}: The process of exchanging a basic and a non-basic variable in the Simplex method to move to a new vertex in the feasible region.<L[2024-08-27-ExSimplex] 9><L[2024-08-27-ExSimplex] 10><L[2024-08-27-ExSimplex] 12><L[2024-08-27-ExSimplex] 13><L[2024-08-27-ExSimplex] 14><L[2024-08-27-ExSimplex] 15>%
\item \href{https://www.math.purdue.edu/~yipn/421/2024-08-27-ExSimplex.pdf#page=26}{\textbf{Dictionary}}: A representation of a linear programming problem where basic variables are expressed as functions of non-basic variables.<L[2024-08-27-ExSimplex] 26><L[2024-08-27-ExSimplex] 27><L[2024-08-27-ExSimplex] 28><L[2024-08-27-ExSimplex] 29>%
\item \href{https://www.math.purdue.edu/~yipn/421/2024-08-27-ExSimplex.pdf#page=32}{\textbf{Auxiliary Problem}}: A modified linear programming problem introduced to find an initial feasible solution when the origin is not feasible in the original problem.<L[2024-08-27-ExSimplex] 32><L[2024-08-27-ExSimplex] 33><L[2024-08-27-ExSimplex] 34><L[2024-08-27-ExSimplex] 35><L[2024-08-27-ExSimplex] 36><L[2024-08-27-ExSimplex] 37><L[2024-08-27-ExSimplex] 38><L[2024-08-27-ExSimplex] 39><L[2024-08-27-ExSimplex] 40><L[2024-08-27-ExSimplex] 41><L[2024-08-27-ExSimplex] 42>%
\end{itemize}%
\begin{itemize}[[leftmargin=*]]%
\item \href{https://www.math.purdue.edu/~yipn/421/2024-09-05-SimplexEfficiency.pdf#page=4}{\textbf{Largest{-}Coefficient Rule}}: A pivot rule in the Simplex method that selects the entering variable with the largest coefficient in the objective function row.<L[2024-09-05-SimplexEfficiency] 4><L[2024-09-05-SimplexEfficiency] 6><L[2024-09-05-SimplexEfficiency] 7><L[2024-09-05-SimplexEfficiency] 8>%
\item \href{https://www.math.purdue.edu/~yipn/421/2024-09-05-SimplexEfficiency.pdf#page=6}{\textbf{Largest{-}Increase Rule}}: A pivot rule in the Simplex method that selects the entering variable that produces the largest increase in the objective function value.<L[2024-09-05-SimplexEfficiency] 6><L[2024-09-05-SimplexEfficiency] 7><L[2024-09-05-SimplexEfficiency] 8>%
\item \href{https://www.math.purdue.edu/~yipn/421/2024-09-05-SimplexEfficiency.pdf#page=2}{\textbf{Klee{-}Minty Example}}: A worst-case example demonstrating that the Simplex method can take an exponential number of iterations.<L[2024-09-05-SimplexEfficiency] 2><L[2024-09-05-SimplexEfficiency] 3><L[2024-09-05-SimplexEfficiency] 4><L[2024-09-05-SimplexEfficiency] 5>%
\item \href{https://www.math.purdue.edu/~yipn/421/2024-09-05-SimplexEfficiency.pdf#page=10}{\textbf{Bland's Rule}}: A pivot rule in the Simplex method designed to prevent cycling.<L[2024-09-05-SimplexEfficiency] 10>%
\item \href{https://www.math.purdue.edu/~yipn/421/2024-09-05-SimplexEfficiency.pdf#page=10}{\textbf{Lexicographic Method}}: A method for choosing the leaving variable in the Simplex method to prevent cycling.<L[2024-09-05-SimplexEfficiency] 10>%
\item \href{https://www.math.purdue.edu/~yipn/421/2024-09-05-SimplexEfficiency.pdf#page=11}{\textbf{Smoothed Complexity}}: A framework for analyzing the complexity of algorithms by considering the average-case performance after small random perturbations to the input.<L[2024-09-05-SimplexEfficiency] 11>%
\item \href{https://www.math.purdue.edu/~yipn/421/2024-09-05-SimplexEfficiency.pdf#page=12}{\textbf{Interior Method}}: A class of linear programming algorithms that traverse the interior of the feasible region, in contrast to the Simplex method which follows the boundary.<L[2024-09-05-SimplexEfficiency] 12>%
\end{itemize}%
\begin{itemize}[[leftmargin=*]]%
\item \href{https://www.math.purdue.edu/~yipn/421/2024-09-12-Duality.pdf#page=1}{\textbf{Primal Problem}}: A linear programming problem in its standard maximization form, aiming to maximize a linear objective function subject to linear inequality constraints and non-negativity constraints.<L[2024-09-12-Duality] 1><L[2024-09-12-Duality] 2><L[2024-09-12-Duality] 3><L[2024-09-12-Duality] 4><L[2024-09-12-Duality] 6><L[2024-09-12-Duality] 7><L[2024-09-12-Duality] 8><L[2024-09-12-Duality] 9><L[2024-09-12-Duality] 10><L[2024-09-12-Duality] 12><L[2024-09-12-Duality] 13><L[2024-09-12-Duality] 14><L[2024-09-12-Duality] 15><L[2024-09-12-Duality] 16><L[2024-09-12-Duality] 17><L[2024-09-12-Duality] 18><L[2024-09-12-Duality] 19><L[2024-09-12-Duality] 20><L[2024-09-12-Duality] 22><L[2024-09-12-Duality] 23>%
\item \href{https://www.math.purdue.edu/~yipn/421/2024-09-12-Duality.pdf#page=1}{\textbf{Dual Problem}}: A linear programming problem derived from the primal problem, aiming to minimize a linear objective function subject to linear inequality constraints and non-negativity constraints. Its solution provides information about the primal problem's solution.<L[2024-09-12-Duality] 1><L[2024-09-12-Duality] 2><L[2024-09-12-Duality] 3><L[2024-09-12-Duality] 4><L[2024-09-12-Duality] 6><L[2024-09-12-Duality] 7><L[2024-09-12-Duality] 8><L[2024-09-12-Duality] 9><L[2024-09-12-Duality] 10><L[2024-09-12-Duality] 11><L[2024-09-12-Duality] 17><L[2024-09-12-Duality] 18><L[2024-09-12-Duality] 19><L[2024-09-12-Duality] 20><L[2024-09-12-Duality] 21><L[2024-09-12-Duality] 22><L[2024-09-12-Duality] 23>%
\item \href{https://www.math.purdue.edu/~yipn/421/2024-09-12-Duality.pdf#page=2}{\textbf{Weak Duality}}: A theorem stating that the objective function value of any feasible solution to the primal problem is always less than or equal to the objective function value of any feasible solution to the dual problem.<L[2024-09-12-Duality] 2>%
\item \href{https://www.math.purdue.edu/~yipn/421/2024-09-12-Duality.pdf#page=6}{\textbf{Negative Transpose Property}}: The property that the coefficient matrix of the dual problem is the negative transpose of the coefficient matrix of the primal problem. This property is preserved throughout the simplex method iterations.<L[2024-09-12-Duality] 6><L[2024-09-12-Duality] 7><L[2024-09-12-Duality] 8><L[2024-09-12-Duality] 9><L[2024-09-12-Duality] 10>%
\item \href{https://www.math.purdue.edu/~yipn/421/2024-09-12-Duality.pdf#page=5}{\textbf{Complementary Variables}}: Pairs of primal and dual variables (and their corresponding slack variables) that exhibit a complementary relationship in the optimal solution; at optimality, at least one variable in each pair must be zero.<L[2024-09-12-Duality] 5><L[2024-09-12-Duality] 21><L[2024-09-12-Duality] 23>%
\item \href{https://www.math.purdue.edu/~yipn/421/2024-09-12-Duality.pdf#page=10}{\textbf{Strong Duality}}: A theorem stating that if either the primal or dual problem has an optimal solution, then so does the other, and the optimal objective function values are equal. The duality gap is zero.<L[2024-09-12-Duality] 10><L[2024-09-12-Duality] 11><L[2024-09-12-Duality] 17><L[2024-09-12-Duality] 18><L[2024-09-12-Duality] 19><L[2024-09-12-Duality] 20>%
\item \href{https://www.math.purdue.edu/~yipn/421/2024-09-12-Duality.pdf#page=21}{\textbf{Complementary Slackness}}: A condition stating that for optimal primal and dual solutions, the product of each primal variable and its corresponding dual slack variable is zero, and the product of each dual variable and its corresponding primal slack variable is zero.<L[2024-09-12-Duality] 21><L[2024-09-12-Duality] 23>%
\end{itemize}%
\begin{itemize}[[leftmargin=*]]%
\item \href{https://www.math.purdue.edu/~yipn/421/2024-09-17-DualityGeneral.pdf#page=1}{\textbf{Primal Problem}}: A linear programming problem in its original or standard form, typically aiming to maximize an objective function subject to constraints.<L[2024-09-17-DualityGeneral] 1><L[2024-09-17-DualityGeneral] 6><L[2024-09-17-DualityGeneral] 7><L[2024-09-17-DualityGeneral] 10><L[2024-09-17-DualityGeneral] 11><L[2024-09-17-DualityGeneral] 12><L[2024-09-17-DualityGeneral] 13><L[2024-09-17-DualityGeneral] 14><L[2024-09-17-DualityGeneral] 15><L[2024-09-17-DualityGeneral] 16><L[2024-09-17-DualityGeneral] 17><L[2024-09-17-DualityGeneral] 18><L[2024-09-17-DualityGeneral] 19><L[2024-09-17-DualityGeneral] 20><L[2024-09-17-DualityGeneral] 21><L[2024-09-17-DualityGeneral] 22>%
\item \href{https://www.math.purdue.edu/~yipn/421/2024-09-17-DualityGeneral.pdf#page=1}{\textbf{Dual Problem}}: A linear programming problem derived from the primal problem, typically aiming to minimize an objective function subject to constraints derived from the primal problem's constraints and objective function.<L[2024-09-17-DualityGeneral] 1><L[2024-09-17-DualityGeneral] 6><L[2024-09-17-DualityGeneral] 7><L[2024-09-17-DualityGeneral] 10><L[2024-09-17-DualityGeneral] 11><L[2024-09-17-DualityGeneral] 12><L[2024-09-17-DualityGeneral] 13><L[2024-09-17-DualityGeneral] 14><L[2024-09-17-DualityGeneral] 15><L[2024-09-17-DualityGeneral] 16><L[2024-09-17-DualityGeneral] 17><L[2024-09-17-DualityGeneral] 18><L[2024-09-17-DualityGeneral] 19><L[2024-09-17-DualityGeneral] 20><L[2024-09-17-DualityGeneral] 21><L[2024-09-17-DualityGeneral] 22>%
\item \href{https://www.math.purdue.edu/~yipn/421/2024-09-17-DualityGeneral.pdf#page=1}{\textbf{Objective Function}}: The function to be maximized or minimized in a linear programming problem. It represents the goal of the optimization.<L[2024-09-17-DualityGeneral] 1><L[2024-09-17-DualityGeneral] 2><L[2024-09-17-DualityGeneral] 3><L[2024-09-17-DualityGeneral] 4><L[2024-09-17-DualityGeneral] 5><L[2024-09-17-DualityGeneral] 6><L[2024-09-17-DualityGeneral] 7><L[2024-09-17-DualityGeneral] 8><L[2024-09-17-DualityGeneral] 9><L[2024-09-17-DualityGeneral] 10><L[2024-09-17-DualityGeneral] 11><L[2024-09-17-DualityGeneral] 12><L[2024-09-17-DualityGeneral] 13><L[2024-09-17-DualityGeneral] 14><L[2024-09-17-DualityGeneral] 15><L[2024-09-17-DualityGeneral] 16><L[2024-09-17-DualityGeneral] 17><L[2024-09-17-DualityGeneral] 18><L[2024-09-17-DualityGeneral] 19><L[2024-09-17-DualityGeneral] 20><L[2024-09-17-DualityGeneral] 21><L[2024-09-17-DualityGeneral] 22>%
\item \href{https://www.math.purdue.edu/~yipn/421/2024-09-17-DualityGeneral.pdf#page=1}{\textbf{Constraints}}: Conditions or restrictions that limit the feasible region of solutions in a linear programming problem. They can be equality constraints (=) or inequality constraints (≤ or ≥).<L[2024-09-17-DualityGeneral] 1><L[2024-09-17-DualityGeneral] 2><L[2024-09-17-DualityGeneral] 3><L[2024-09-17-DualityGeneral] 4><L[2024-09-17-DualityGeneral] 5><L[2024-09-17-DualityGeneral] 6><L[2024-09-17-DualityGeneral] 7><L[2024-09-17-DualityGeneral] 8><L[2024-09-17-DualityGeneral] 9><L[2024-09-17-DualityGeneral] 10><L[2024-09-17-DualityGeneral] 11><L[2024-09-17-DualityGeneral] 12><L[2024-09-17-DualityGeneral] 13><L[2024-09-17-DualityGeneral] 14><L[2024-09-17-DualityGeneral] 15><L[2024-09-17-DualityGeneral] 16><L[2024-09-17-DualityGeneral] 17><L[2024-09-17-DualityGeneral] 18><L[2024-09-17-DualityGeneral] 19><L[2024-09-17-DualityGeneral] 20><L[2024-09-17-DualityGeneral] 21><L[2024-09-17-DualityGeneral] 22>%
\item \href{https://www.math.purdue.edu/~yipn/421/2024-09-17-DualityGeneral.pdf#page=2}{\textbf{Lagrange Multiplier}}: A variable introduced in the Lagrangian function to incorporate constraints into the optimization problem. They represent the sensitivity of the objective function to changes in the constraints.<L[2024-09-17-DualityGeneral] 2><L[2024-09-17-DualityGeneral] 7><L[2024-09-17-DualityGeneral] 8><L[2024-09-17-DualityGeneral] 9><L[2024-09-17-DualityGeneral] 10><L[2024-09-17-DualityGeneral] 11><L[2024-09-17-DualityGeneral] 14><L[2024-09-17-DualityGeneral] 15><L[2024-09-17-DualityGeneral] 16><L[2024-09-17-DualityGeneral] 17><L[2024-09-17-DualityGeneral] 19><L[2024-09-17-DualityGeneral] 20>%
\item \href{https://www.math.purdue.edu/~yipn/421/2024-09-17-DualityGeneral.pdf#page=2}{\textbf{Lagrangian Function}}: A function that combines the objective function and constraints of a linear programming problem using Lagrange multipliers. Its maximization/minimization helps solve the constrained optimization problem.<L[2024-09-17-DualityGeneral] 2><L[2024-09-17-DualityGeneral] 7><L[2024-09-17-DualityGeneral] 8><L[2024-09-17-DualityGeneral] 9><L[2024-09-17-DualityGeneral] 10><L[2024-09-17-DualityGeneral] 11><L[2024-09-17-DualityGeneral] 14><L[2024-09-17-DualityGeneral] 15><L[2024-09-17-DualityGeneral] 16><L[2024-09-17-DualityGeneral] 17><L[2024-09-17-DualityGeneral] 19><L[2024-09-17-DualityGeneral] 20>%
\item \href{https://www.math.purdue.edu/~yipn/421/2024-09-17-DualityGeneral.pdf#page=5}{\textbf{Weak Duality}}: A theorem stating that the optimal value of the primal problem is always less than or equal to the optimal value of the dual problem.<L[2024-09-17-DualityGeneral] 5><L[2024-09-17-DualityGeneral] 6><L[2024-09-17-DualityGeneral] 9><L[2024-09-17-DualityGeneral] 10><L[2024-09-17-DualityGeneral] 18>%
\item \href{https://www.math.purdue.edu/~yipn/421/2024-09-17-DualityGeneral.pdf#page=17}{\textbf{Strong Duality}}: A theorem stating that if the primal problem has a finite optimal solution, then the dual problem also has a finite optimal solution, and their optimal values are equal.<L[2024-09-17-DualityGeneral] 17><L[2024-09-17-DualityGeneral] 21><L[2024-09-17-DualityGeneral] 22>%
\item \href{https://www.math.purdue.edu/~yipn/421/2024-09-17-DualityGeneral.pdf#page=22}{\textbf{Complementary Slackness}}: A condition that holds at the optimal solution of a linear program and its dual. It states that for each constraint, either the constraint is binding (holds with equality) or the corresponding dual variable is zero.<L[2024-09-17-DualityGeneral] 22>%
\end{itemize}%
\begin{itemize}[[leftmargin=*]]%
\item \href{https://www.math.purdue.edu/~yipn/421/2024-09-19-DualityExamples.pdf#page=1}{\textbf{Primal Problem}}: The original linear programming problem, typically formulated as a maximization problem. <L[2024-09-19-DualityExamples] 1>%
\item \href{https://www.math.purdue.edu/~yipn/421/2024-09-19-DualityExamples.pdf#page=1}{\textbf{Dual Problem}}: The corresponding minimization problem derived from the primal problem, offering an alternative perspective on the same optimization problem. <L[2024-09-19-DualityExamples] 1>%
\item \href{https://www.math.purdue.edu/~yipn/421/2024-09-19-DualityExamples.pdf#page=2}{\textbf{Diet Problem}}: A classic linear programming problem where the goal is to minimize the cost of a diet while meeting minimum nutritional requirements. <L[2024-09-19-DualityExamples] 2><L[2024-09-19-DualityExamples] 3><L[2024-09-19-DualityExamples] 4><L[2024-09-19-DualityExamples] 5><L[2024-09-19-DualityExamples] 6><L[2024-09-19-DualityExamples] 7><L[2024-09-19-DualityExamples] 8>%
\item \href{https://www.math.purdue.edu/~yipn/421/2024-09-19-DualityExamples.pdf#page=2}{\textbf{Nutrients}}: The essential components (e.g., vitamins, minerals) that need to be included in a diet. <L[2024-09-19-DualityExamples] 2><L[2024-09-19-DualityExamples] 3><L[2024-09-19-DualityExamples] 4><L[2024-09-19-DualityExamples] 6><L[2024-09-19-DualityExamples] 7>%
\item \href{https://www.math.purdue.edu/~yipn/421/2024-09-19-DualityExamples.pdf#page=2}{\textbf{Foods}}: The different food items available to be included in the diet. <L[2024-09-19-DualityExamples] 2><L[2024-09-19-DualityExamples] 3><L[2024-09-19-DualityExamples] 4><L[2024-09-19-DualityExamples] 6><L[2024-09-19-DualityExamples] 7>%
\item \href{https://www.math.purdue.edu/~yipn/421/2024-09-19-DualityExamples.pdf#page=1}{\textbf{Constraint Matrix (A)}}: A matrix representing the relationships between the decision variables and the constraints in a linear programming problem. <L[2024-09-19-DualityExamples] 1>%
\item \href{https://www.math.purdue.edu/~yipn/421/2024-09-19-DualityExamples.pdf#page=1}{\textbf{Objective Function}}: The function to be maximized or minimized in a linear programming problem. <L[2024-09-19-DualityExamples] 1><L[2024-09-19-DualityExamples] 5><L[2024-09-19-DualityExamples] 8>%
\item \href{https://www.math.purdue.edu/~yipn/421/2024-09-19-DualityExamples.pdf#page=1}{\textbf{Decision Variables}}: The variables representing the quantities to be determined in a linear programming problem (e.g., the amount of each food to consume). <L[2024-09-19-DualityExamples] 1><L[2024-09-19-DualityExamples] 2><L[2024-09-19-DualityExamples] 5><L[2024-09-19-DualityExamples] 8>%
\end{itemize}%
\begin{itemize}[[leftmargin=*]]%
\item \href{https://www.math.purdue.edu/~yipn/421/2024-09-24-SimplexMatrix.pdf#page=4}{\textbf{Basic Variables}}: Variables in the simplex method that are set to non-zero values to determine a basic feasible solution.<L[2024-09-24-SimplexMatrix] 4><L[2024-09-24-SimplexMatrix] 5><L[2024-09-24-SimplexMatrix] 6><L[2024-09-24-SimplexMatrix] 7><L[2024-09-24-SimplexMatrix] 8><L[2024-09-24-SimplexMatrix] 9><L[2024-09-24-SimplexMatrix] 10><L[2024-09-24-SimplexMatrix] 11><L[2024-09-24-SimplexMatrix] 12><L[2024-09-24-SimplexMatrix] 13><L[2024-09-24-SimplexMatrix] 14><L[2024-09-24-SimplexMatrix] 15><L[2024-09-24-SimplexMatrix] 17><L[2024-09-24-SimplexMatrix] 18><L[2024-09-24-SimplexMatrix] 19><L[2024-09-24-SimplexMatrix] 20>%
\item \href{https://www.math.purdue.edu/~yipn/421/2024-09-24-SimplexMatrix.pdf#page=4}{\textbf{Non{-}Basic Variables}}: Variables in the simplex method that are set to zero to determine a basic feasible solution.<L[2024-09-24-SimplexMatrix] 4><L[2024-09-24-SimplexMatrix] 5><L[2024-09-24-SimplexMatrix] 6><L[2024-09-24-SimplexMatrix] 7><L[2024-09-24-SimplexMatrix] 8><L[2024-09-24-SimplexMatrix] 9><L[2024-09-24-SimplexMatrix] 10><L[2024-09-24-SimplexMatrix] 11><L[2024-09-24-SimplexMatrix] 12><L[2024-09-24-SimplexMatrix] 13><L[2024-09-24-SimplexMatrix] 14><L[2024-09-24-SimplexMatrix] 15><L[2024-09-24-SimplexMatrix] 17><L[2024-09-24-SimplexMatrix] 18><L[2024-09-24-SimplexMatrix] 19><L[2024-09-24-SimplexMatrix] 20>%
\item \href{https://www.math.purdue.edu/~yipn/421/2024-09-24-SimplexMatrix.pdf#page=1}{\textbf{Slack Variables}}: Variables added to inequality constraints to transform them into equality constraints.<L[2024-09-24-SimplexMatrix] 1><L[2024-09-24-SimplexMatrix] 2><L[2024-09-24-SimplexMatrix] 3><L[2024-09-24-SimplexMatrix] 5><L[2024-09-24-SimplexMatrix] 6><L[2024-09-24-SimplexMatrix] 7>%
\item \href{https://www.math.purdue.edu/~yipn/421/2024-09-24-SimplexMatrix.pdf#page=1}{\textbf{Simplex Method}}: An iterative algorithm used to solve linear programming problems by moving from one basic feasible solution to another until the optimal solution is found.<L[2024-09-24-SimplexMatrix] 1><L[2024-09-24-SimplexMatrix] 2><L[2024-09-24-SimplexMatrix] 3><L[2024-09-24-SimplexMatrix] 4><L[2024-09-24-SimplexMatrix] 5><L[2024-09-24-SimplexMatrix] 6><L[2024-09-24-SimplexMatrix] 7><L[2024-09-24-SimplexMatrix] 8><L[2024-09-24-SimplexMatrix] 9><L[2024-09-24-SimplexMatrix] 10><L[2024-09-24-SimplexMatrix] 11><L[2024-09-24-SimplexMatrix] 12><L[2024-09-24-SimplexMatrix] 13><L[2024-09-24-SimplexMatrix] 14><L[2024-09-24-SimplexMatrix] 15><L[2024-09-24-SimplexMatrix] 17><L[2024-09-24-SimplexMatrix] 18><L[2024-09-24-SimplexMatrix] 19><L[2024-09-24-SimplexMatrix] 20>%
\item \href{https://www.math.purdue.edu/~yipn/421/2024-09-24-SimplexMatrix.pdf#page=1}{\textbf{Objective Function}}: The function to be maximized or minimized in a linear programming problem.<L[2024-09-24-SimplexMatrix] 1><L[2024-09-24-SimplexMatrix] 2><L[2024-09-24-SimplexMatrix] 3><L[2024-09-24-SimplexMatrix] 4><L[2024-09-24-SimplexMatrix] 8><L[2024-09-24-SimplexMatrix] 9><L[2024-09-24-SimplexMatrix] 10><L[2024-09-24-SimplexMatrix] 11><L[2024-09-24-SimplexMatrix] 12><L[2024-09-24-SimplexMatrix] 13><L[2024-09-24-SimplexMatrix] 15><L[2024-09-24-SimplexMatrix] 17><L[2024-09-24-SimplexMatrix] 18><L[2024-09-24-SimplexMatrix] 19><L[2024-09-24-SimplexMatrix] 20><L[2024-09-24-SimplexMatrix] 21>%
\item \href{https://www.math.purdue.edu/~yipn/421/2024-09-24-SimplexMatrix.pdf#page=1}{\textbf{Constraints}}: Conditions that restrict the values of the variables in a linear programming problem.<L[2024-09-24-SimplexMatrix] 1><L[2024-09-24-SimplexMatrix] 2><L[2024-09-24-SimplexMatrix] 3><L[2024-09-24-SimplexMatrix] 5><L[2024-09-24-SimplexMatrix] 6><L[2024-09-24-SimplexMatrix] 7><L[2024-09-24-SimplexMatrix] 9><L[2024-09-24-SimplexMatrix] 12><L[2024-09-24-SimplexMatrix] 13><L[2024-09-24-SimplexMatrix] 15><L[2024-09-24-SimplexMatrix] 17><L[2024-09-24-SimplexMatrix] 18><L[2024-09-24-SimplexMatrix] 19><L[2024-09-24-SimplexMatrix] 20><L[2024-09-24-SimplexMatrix] 21>%
\item \href{https://www.math.purdue.edu/~yipn/421/2024-09-24-SimplexMatrix.pdf#page=1}{\textbf{Linear Programming}}: A mathematical method for determining the optimal allocation of resources subject to constraints.<L[2024-09-24-SimplexMatrix] 1><L[2024-09-24-SimplexMatrix] 2><L[2024-09-24-SimplexMatrix] 3><L[2024-09-24-SimplexMatrix] 15><L[2024-09-24-SimplexMatrix] 21>%
\item \href{https://www.math.purdue.edu/~yipn/421/2024-09-24-SimplexMatrix.pdf#page=14}{\textbf{Complementary Slackness}}: A condition that relates the primal and dual variables in a linear programming problem at optimality.<L[2024-09-24-SimplexMatrix] 14>%
\end{itemize}%
\begin{itemize}[[leftmargin=*]]%
\item \href{https://www.math.purdue.edu/~yipn/421/2024-09-26-Sensitivity.pdf#page=1}{\textbf{Objective Function}}: A mathematical expression that defines the quantity to be maximized or minimized in a linear programming problem.<L[2024-09-26-Sensitivity] 1><L[2024-09-26-Sensitivity] 4><L[2024-09-26-Sensitivity] 8>%
\item \href{https://www.math.purdue.edu/~yipn/421/2024-09-26-Sensitivity.pdf#page=1}{\textbf{Constraints}}: Limitations or restrictions on the values of the variables in a linear programming problem, expressed as inequalities or equations.<L[2024-09-26-Sensitivity] 1><L[2024-09-26-Sensitivity] 4><L[2024-09-26-Sensitivity] 8>%
\item \href{https://www.math.purdue.edu/~yipn/421/2024-09-26-Sensitivity.pdf#page=2}{\textbf{Optimal Solution}}: The set of variable values that satisfies all constraints and yields the best possible value of the objective function.<L[2024-09-26-Sensitivity] 2><L[2024-09-26-Sensitivity] 3><L[2024-09-26-Sensitivity] 5><L[2024-09-26-Sensitivity] 6><L[2024-09-26-Sensitivity] 9><L[2024-09-26-Sensitivity] 10>%
\item \href{https://www.math.purdue.edu/~yipn/421/2024-09-26-Sensitivity.pdf#page=1}{\textbf{Sensitivity Analysis}}: The study of how changes in the parameters of a linear programming problem (objective function coefficients or constraint values) affect the optimal solution.<L[2024-09-26-Sensitivity] 1><L[2024-09-26-Sensitivity] 4><L[2024-09-26-Sensitivity] 8>%
\item \href{https://www.math.purdue.edu/~yipn/421/2024-09-26-Sensitivity.pdf#page=4}{\textbf{Parametric Analysis}}: A type of sensitivity analysis where one or more parameters are systematically varied to observe their impact on the optimal solution.<L[2024-09-26-Sensitivity] 4><L[2024-09-26-Sensitivity] 8>%
\item \href{https://www.math.purdue.edu/~yipn/421/2024-09-26-Sensitivity.pdf#page=2}{\textbf{Basic Variables}}: Variables that have non-zero values in the optimal solution of a linear programming problem, typically corresponding to the columns of the basis matrix.<L[2024-09-26-Sensitivity] 2><L[2024-09-26-Sensitivity] 3><L[2024-09-26-Sensitivity] 5><L[2024-09-26-Sensitivity] 6><L[2024-09-26-Sensitivity] 9><L[2024-09-26-Sensitivity] 10><L[2024-09-26-Sensitivity] 7><L[2024-09-26-Sensitivity] 11>%
\item \href{https://www.math.purdue.edu/~yipn/421/2024-09-26-Sensitivity.pdf#page=2}{\textbf{Non{-}basic Variables}}: Variables that have zero values in the optimal solution of a linear programming problem.<L[2024-09-26-Sensitivity] 2><L[2024-09-26-Sensitivity] 3><L[2024-09-26-Sensitivity] 5><L[2024-09-26-Sensitivity] 6><L[2024-09-26-Sensitivity] 9><L[2024-09-26-Sensitivity] 10><L[2024-09-26-Sensitivity] 7><L[2024-09-26-Sensitivity] 11>%
\item \href{https://www.math.purdue.edu/~yipn/421/2024-09-26-Sensitivity.pdf#page=7}{\textbf{Reduced Costs}}: Values that indicate how much the objective function value would change if a non-basic variable were to become basic.<L[2024-09-26-Sensitivity] 7>%
\item \href{https://www.math.purdue.edu/~yipn/421/2024-09-26-Sensitivity.pdf#page=8}{\textbf{Right{-}Hand Side (RHS) Values}}: The constants on the right-hand side of the constraint inequalities in a linear programming problem.<L[2024-09-26-Sensitivity] 8>%
\item \href{https://www.math.purdue.edu/~yipn/421/2024-09-26-Sensitivity.pdf#page=7}{\textbf{Basis Matrix (B)}}: A square matrix formed by the columns of the constraint matrix corresponding to the basic variables.<L[2024-09-26-Sensitivity] 7><L[2024-09-26-Sensitivity] 11>%
\end{itemize}%
\begin{itemize}[[leftmargin=*]]%
\item \href{https://www.math.purdue.edu/~yipn/421/2024-10-01-DualGeneralLP.pdf#page=1}{\textbf{Primal Problem (P)}}: The original linear programming problem, typically formulated as a maximization problem.<L[2024-10-01-DualGeneralLP] 1>%
\item \href{https://www.math.purdue.edu/~yipn/421/2024-10-01-DualGeneralLP.pdf#page=1}{\textbf{Dual Problem (D)}}: A linear programming problem derived from the primal problem, typically formulated as a minimization problem.<L[2024-10-01-DualGeneralLP] 1><L[2024-10-01-DualGeneralLP] 3><L[2024-10-01-DualGeneralLP] 7>%
\item \href{https://www.math.purdue.edu/~yipn/421/2024-10-01-DualGeneralLP.pdf#page=1}{\textbf{Lagrange Multiplier}}: A variable used in the Lagrangian function to incorporate constraints into the objective function.<L[2024-10-01-DualGeneralLP] 1><L[2024-10-01-DualGeneralLP] 4>%
\item \href{https://www.math.purdue.edu/~yipn/421/2024-10-01-DualGeneralLP.pdf#page=4}{\textbf{Lagrangian Function}}: A function that combines the objective function and constraints of a primal problem using Lagrange multipliers.<L[2024-10-01-DualGeneralLP] 4>%
\item \href{https://www.math.purdue.edu/~yipn/421/2024-10-01-DualGeneralLP.pdf#page=5}{\textbf{Weak Duality}}: A theorem stating that the optimal value of the primal problem is always less than or equal to the optimal value of the dual problem.<L[2024-10-01-DualGeneralLP] 5>%
\end{itemize}%
\begin{itemize}[[leftmargin=*]]%
\item \href{https://www.math.purdue.edu/~yipn/421/2024-10-01-SensitivityDuality.pdf#page=1}{\textbf{Primal Problem}}: A linear programming problem in its standard maximization form.<L[2024-10-01-SensitivityDuality] 1>%
\item \href{https://www.math.purdue.edu/~yipn/421/2024-10-01-SensitivityDuality.pdf#page=2}{\textbf{Dual Problem}}: The linear programming problem associated with a given primal problem, formulated to minimize the objective function subject to constraints derived from the primal problem.<L[2024-10-01-SensitivityDuality] 2>%
\item \href{https://www.math.purdue.edu/~yipn/421/2024-10-01-SensitivityDuality.pdf#page=2}{\textbf{Strong Duality}}: A fundamental theorem in linear programming stating that if both the primal and dual problems have feasible solutions, then their optimal objective function values are equal.<L[2024-10-01-SensitivityDuality] 2>%
\item \href{https://www.math.purdue.edu/~yipn/421/2024-10-01-SensitivityDuality.pdf#page=3}{\textbf{Complementary Slackness}}: A condition that must hold at the optimal solution of a linear programming problem and its dual, relating the primal variables and dual slack variables.<L[2024-10-01-SensitivityDuality] 3>%
\item \href{https://www.math.purdue.edu/~yipn/421/2024-10-01-SensitivityDuality.pdf#page=4}{\textbf{Basis (B)}}: A set of linearly independent columns from the constraint matrix of a linear programming problem, used in the simplex method to represent the basic variables.<L[2024-10-01-SensitivityDuality] 4><L[2024-10-01-SensitivityDuality] 6><L[2024-10-01-SensitivityDuality] 11><L[2024-10-01-SensitivityDuality] 12><L[2024-10-01-SensitivityDuality] 13><L[2024-10-01-SensitivityDuality] 14><L[2024-10-01-SensitivityDuality] 15>%
\item \href{https://www.math.purdue.edu/~yipn/421/2024-10-01-SensitivityDuality.pdf#page=4}{\textbf{Non{-}basis (N)}}: The set of columns from the constraint matrix not included in the basis.<L[2024-10-01-SensitivityDuality] 4><L[2024-10-01-SensitivityDuality] 6><L[2024-10-01-SensitivityDuality] 11><L[2024-10-01-SensitivityDuality] 12><L[2024-10-01-SensitivityDuality] 15>%
\item \href{https://www.math.purdue.edu/~yipn/421/2024-10-01-SensitivityDuality.pdf#page=4}{\textbf{Simplex Method}}: An iterative algorithm for solving linear programming problems by systematically moving from one feasible solution to another, improving the objective function value at each step.<L[2024-10-01-SensitivityDuality] 4><L[2024-10-01-SensitivityDuality] 5><L[2024-10-01-SensitivityDuality] 6>%
\item \href{https://www.math.purdue.edu/~yipn/421/2024-10-01-SensitivityDuality.pdf#page=1}{\textbf{Sensitivity Analysis}}: The study of how changes in the parameters of a linear programming problem (e.g., objective function coefficients, constraint coefficients, right-hand side values) affect the optimal solution.<L[2024-10-01-SensitivityDuality] 1>%
\item \href{https://www.math.purdue.edu/~yipn/421/2024-10-01-SensitivityDuality.pdf#page=10}{\textbf{Shadow Price}}: The rate of change in the optimal objective function value with respect to a unit change in the right-hand side of a constraint. It is given by the optimal dual variable corresponding to that constraint.<L[2024-10-01-SensitivityDuality] 10>%
\item \href{https://www.math.purdue.edu/~yipn/421/2024-10-01-SensitivityDuality.pdf#page=11}{\textbf{Reduced Costs (z<sub>N</sub>)}}: In the simplex method, the values indicating how much the objective function would change if a non-basic variable were to enter the basis.<L[2024-10-01-SensitivityDuality] 11><L[2024-10-01-SensitivityDuality] 12><L[2024-10-01-SensitivityDuality] 15>%
\item \href{https://www.math.purdue.edu/~yipn/421/2024-10-01-SensitivityDuality.pdf#page=13}{\textbf{Feasibility}}: The condition that a solution satisfies all the constraints of a linear programming problem.<L[2024-10-01-SensitivityDuality] 13><L[2024-10-01-SensitivityDuality] 14>%
\end{itemize}%
\begin{itemize}[[leftmargin=*]]%
\item \href{https://www.math.purdue.edu/~yipn/421/2024-10-03-DualSimplexMatrix.pdf#page=2}{\textbf{Primal Problem}}: A linear programming problem in its original, standard form. <L[2024-10-03-DualSimplexMatrix] 2><L[2024-10-03-DualSimplexMatrix] 3><L[2024-10-03-DualSimplexMatrix] 5><L[2024-10-03-DualSimplexMatrix] 6>%
\item \href{https://www.math.purdue.edu/~yipn/421/2024-10-03-DualSimplexMatrix.pdf#page=2}{\textbf{Dual Problem}}: A linear programming problem derived from the primal problem, offering an alternative perspective for finding the optimal solution. <L[2024-10-03-DualSimplexMatrix] 2><L[2024-10-03-DualSimplexMatrix] 3><L[2024-10-03-DualSimplexMatrix] 5><L[2024-10-03-DualSimplexMatrix] 6>%
\item \href{https://www.math.purdue.edu/~yipn/421/2024-10-03-DualSimplexMatrix.pdf#page=2}{\textbf{Dual Simplex Method}}: An algorithm for solving linear programming problems by iteratively improving the dual feasible solution until the primal feasible solution is reached. <L[2024-10-03-DualSimplexMatrix] 2><L[2024-10-03-DualSimplexMatrix] 3><L[2024-10-03-DualSimplexMatrix] 4><L[2024-10-03-DualSimplexMatrix] 5><L[2024-10-03-DualSimplexMatrix] 6>%
\item \href{https://www.math.purdue.edu/~yipn/421/2024-10-03-DualSimplexMatrix.pdf#page=5}{\textbf{Dual Based Phase I Algorithm}}: A method to find an initial feasible solution for the dual problem, which is then used to solve the primal problem. <L[2024-10-03-DualSimplexMatrix] 5><L[2024-10-03-DualSimplexMatrix] 6>%
\item \href{https://www.math.purdue.edu/~yipn/421/2024-10-03-DualSimplexMatrix.pdf#page=7}{\textbf{Simplex Method in Matrix Form}}: A representation of the simplex algorithm using matrices and vectors, facilitating efficient computation. <L[2024-10-03-DualSimplexMatrix] 7><L[2024-10-03-DualSimplexMatrix] 8><L[2024-10-03-DualSimplexMatrix] 9><L[2024-10-03-DualSimplexMatrix] 10><L[2024-10-03-DualSimplexMatrix] 11><L[2024-10-03-DualSimplexMatrix] 12><L[2024-10-03-DualSimplexMatrix] 13><L[2024-10-03-DualSimplexMatrix] 14><L[2024-10-03-DualSimplexMatrix] 15>%
\item \href{https://www.math.purdue.edu/~yipn/421/2024-10-03-DualSimplexMatrix.pdf#page=7}{\textbf{Basic Variables}}: Variables that form a basis for the solution space in the simplex method. <L[2024-10-03-DualSimplexMatrix] 7><L[2024-10-03-DualSimplexMatrix] 8><L[2024-10-03-DualSimplexMatrix] 9><L[2024-10-03-DualSimplexMatrix] 10><L[2024-10-03-DualSimplexMatrix] 11><L[2024-10-03-DualSimplexMatrix] 12><L[2024-10-03-DualSimplexMatrix] 13><L[2024-10-03-DualSimplexMatrix] 14><L[2024-10-03-DualSimplexMatrix] 15>%
\item \href{https://www.math.purdue.edu/~yipn/421/2024-10-03-DualSimplexMatrix.pdf#page=7}{\textbf{Non{-}basic Variables}}: Variables that are not in the basis in the simplex method. <L[2024-10-03-DualSimplexMatrix] 7><L[2024-10-03-DualSimplexMatrix] 8><L[2024-10-03-DualSimplexMatrix] 9><L[2024-10-03-DualSimplexMatrix] 10><L[2024-10-03-DualSimplexMatrix] 11><L[2024-10-03-DualSimplexMatrix] 12><L[2024-10-03-DualSimplexMatrix] 13><L[2024-10-03-DualSimplexMatrix] 14><L[2024-10-03-DualSimplexMatrix] 15>%
\item \href{https://www.math.purdue.edu/~yipn/421/2024-10-03-DualSimplexMatrix.pdf#page=7}{\textbf{Reduced Costs}}: The coefficients of the non-basic variables in the objective function after the basic variables have been expressed in terms of the non-basic variables. <L[2024-10-03-DualSimplexMatrix] 7><L[2024-10-03-DualSimplexMatrix] 8><L[2024-10-03-DualSimplexMatrix] 9><L[2024-10-03-DualSimplexMatrix] 10><L[2024-10-03-DualSimplexMatrix] 11><L[2024-10-03-DualSimplexMatrix] 12><L[2024-10-03-DualSimplexMatrix] 13><L[2024-10-03-DualSimplexMatrix] 14><L[2024-10-03-DualSimplexMatrix] 15>%
\item \href{https://www.math.purdue.edu/~yipn/421/2024-10-03-DualSimplexMatrix.pdf#page=9}{\textbf{Entering Variable}}: The non-basic variable selected to enter the basis in the next iteration of the simplex method. <L[2024-10-03-DualSimplexMatrix] 9><L[2024-10-03-DualSimplexMatrix] 10><L[2024-10-03-DualSimplexMatrix] 11><L[2024-10-03-DualSimplexMatrix] 12><L[2024-10-03-DualSimplexMatrix] 13><L[2024-10-03-DualSimplexMatrix] 14><L[2024-10-03-DualSimplexMatrix] 15>%
\item \href{https://www.math.purdue.edu/~yipn/421/2024-10-03-DualSimplexMatrix.pdf#page=10}{\textbf{Leaving Variable}}: The basic variable selected to leave the basis in the next iteration of the simplex method. <L[2024-10-03-DualSimplexMatrix] 10><L[2024-10-03-DualSimplexMatrix] 11><L[2024-10-03-DualSimplexMatrix] 12><L[2024-10-03-DualSimplexMatrix] 13><L[2024-10-03-DualSimplexMatrix] 14><L[2024-10-03-DualSimplexMatrix] 15>%
\item \href{https://www.math.purdue.edu/~yipn/421/2024-10-03-DualSimplexMatrix.pdf#page=15}{\textbf{Elementary Matrices}}: Matrices used to represent the basis changes in the simplex method. <L[2024-10-03-DualSimplexMatrix] 15>%
\end{itemize}%
\begin{itemize}[[leftmargin=*]]%
\item \href{https://www.math.purdue.edu/~yipn/421/2024-10-22-Convexity.pdf#page=1}{\textbf{Convex Set}}: A set $S$ such that for any two points $z_1, z_2 \in S$, the line segment connecting them is also contained in $S$, i.e., $\lambda z_1 + (1-\lambda)z_2 \in S$ for all $0 \le \lambda \le 1$.<L[2024-10-22-Convexity] 1>%
\item \href{https://www.math.purdue.edu/~yipn/421/2024-10-22-Convexity.pdf#page=1}{\textbf{Convex Combination}}: A weighted sum of points $z_1, z_2, \dots, z_n$ of the form $\sum_{i=1}^n \lambda_i z_i$, where $\lambda_i \ge 0$ for all $i$ and $\sum_{i=1}^n \lambda_i = 1$.<L[2024-10-22-Convexity] 1>%
\item \href{https://www.math.purdue.edu/~yipn/421/2024-10-22-Convexity.pdf#page=5}{\textbf{Convex Hull}}: The smallest convex set containing a given set $S$, denoted as Conv($S$). It is the intersection of all convex sets containing $S$.<L[2024-10-22-Convexity] 5>%
\item \href{https://www.math.purdue.edu/~yipn/421/2024-10-22-Convexity.pdf#page=8}{\textbf{Caratheodory Theorem}}: If $z$ is in the convex hull of a set $S$ in $\mathbb{R}^m$, then $z$ can be written as a convex combination of at most $m+1$ points from $S$.<L[2024-10-22-Convexity] 8><L[2024-10-22-Convexity] 9>%
\item \href{https://www.math.purdue.edu/~yipn/421/2024-10-22-Convexity.pdf#page=10}{\textbf{Separation Theorem}}: If $P$ and $\bar{P}$ are two disjoint nonempty polyhedra in $\mathbb{R}^n$, then there exist disjoint half-spaces $H$ and $\bar{H}$ such that $P \subset H$ and $\bar{P} \subset \bar{H}$.<L[2024-10-22-Convexity] 10><L[2024-10-22-Convexity] 11><L[2024-10-22-Convexity] 12><L[2024-10-22-Convexity] 13><L[2024-10-22-Convexity] 14><L[2024-10-22-Convexity] 15>%
\item \href{https://www.math.purdue.edu/~yipn/421/2024-10-22-Convexity.pdf#page=12}{\textbf{Farkas' Lemma}}: The system of inequalities $Ax \le b$ has no solution if and only if there exists a vector $y \ge 0$ such that $y^T A = 0$ and $y^T b < 0$.<L[2024-10-22-Convexity] 12><L[2024-10-22-Convexity] 13><L[2024-10-22-Convexity] 14><L[2024-10-22-Convexity] 15>%
\end{itemize}%
\begin{itemize}[[leftmargin=*]]%
\item \href{https://www.math.purdue.edu/~yipn/421/2024-10-22-Regression.pdf#page=4}{\textbf{L² Regression (Least Square)}}: A regression technique that minimizes the sum of squared differences between observed and predicted values. <L[2024-10-22-Regression] 4>%
\item \href{https://www.math.purdue.edu/~yipn/421/2024-10-22-Regression.pdf#page=5}{\textbf{L¹ Regression}}: A regression technique that minimizes the sum of absolute differences between observed and predicted values. <L[2024-10-22-Regression] 5>%
\item \href{https://www.math.purdue.edu/~yipn/421/2024-10-22-Regression.pdf#page=6}{\textbf{L∞ Regression}}: A regression technique that minimizes the maximum absolute difference between observed and predicted values. <L[2024-10-22-Regression] 6>%
\item \href{https://www.math.purdue.edu/~yipn/421/2024-10-22-Regression.pdf#page=8}{\textbf{Linear Separability}}: The property of a dataset where a hyperplane can perfectly separate data points into different classes. <L[2024-10-22-Regression] 8>%
\item \href{https://www.math.purdue.edu/~yipn/421/2024-10-22-Regression.pdf#page=9}{\textbf{Maximum Margin Classifier}}: A classifier that aims to maximize the distance between the separating hyperplane and the nearest data points. <L[2024-10-22-Regression] 9>%
\item \href{https://www.math.purdue.edu/~yipn/421/2024-10-22-Regression.pdf#page=10}{\textbf{Support Vector Machine (SVM)}}: A classification algorithm that finds the optimal hyperplane that maximizes the margin between classes. <L[2024-10-22-Regression] 10>%
\item \href{https://www.math.purdue.edu/~yipn/421/2024-10-22-Regression.pdf#page=10}{\textbf{Support Vectors}}: The data points closest to the hyperplane in an SVM, crucial for defining the hyperplane. <L[2024-10-22-Regression] 10><L[2024-10-22-Regression] 17>%
\item \href{https://www.math.purdue.edu/~yipn/421/2024-10-22-Regression.pdf#page=11}{\textbf{Lagrangian Formulation of SVM}}: A method to reformulate the SVM optimization problem using Lagrange multipliers, leading to a dual problem. <L[2024-10-22-Regression] 11>%
\item \href{https://www.math.purdue.edu/~yipn/421/2024-10-22-Regression.pdf#page=12}{\textbf{Geometric Interpretation of SVM}}: Understanding SVM through the geometric concept of maximizing the margin between classes. <L[2024-10-22-Regression] 12><L[2024-10-22-Regression] 17>%
\item \href{https://www.math.purdue.edu/~yipn/421/2024-10-22-Regression.pdf#page=13}{\textbf{SVM Optimization as a QP Problem}}: Formulating the SVM problem as a quadratic programming problem to find the optimal hyperplane. <L[2024-10-22-Regression] 13><L[2024-10-22-Regression] 14><L[2024-10-22-Regression] 19><L[2024-10-22-Regression] 20>%
\item \href{https://www.math.purdue.edu/~yipn/421/2024-10-22-Regression.pdf#page=15}{\textbf{Maximum Inscribed Circle in a Convex Polyhedron}}: A geometric problem of finding the largest circle that fits inside a convex polyhedron. <L[2024-10-22-Regression] 15><L[2024-10-22-Regression] 16><L[2024-10-22-Regression] 21>%
\item \href{https://www.math.purdue.edu/~yipn/421/2024-10-22-Regression.pdf#page=18}{\textbf{Distance of a Point from a Hyperplane}}: The perpendicular distance between a point and a hyperplane, calculated using a specific formula. <L[2024-10-22-Regression] 18>%
\item \href{https://www.math.purdue.edu/~yipn/421/2024-10-22-Regression.pdf#page=22}{\textbf{Smallest Ball Problem}}: Finding the smallest-radius sphere that encloses a given set of points in a d-dimensional space. <L[2024-10-22-Regression] 22>%
\end{itemize}%
\begin{itemize}[[leftmargin=*]]%
\item \href{https://www.math.purdue.edu/~yipn/421/2024-10-28-ApplsLinProg.pdf#page=1}{\textbf{Monthly Demands}}: The predicted sales of ice cream for each month of the year, used as input for the production planning problem.<L[2024-10-28-ApplsLinProg] 1><L[2024-10-28-ApplsLinProg] 2>%
\item \href{https://www.math.purdue.edu/~yipn/421/2024-10-28-ApplsLinProg.pdf#page=2}{\textbf{Average Demand}}: The average monthly sales of ice cream over the entire year, used as a reference point in production planning.<L[2024-10-28-ApplsLinProg] 2>%
\item \href{https://www.math.purdue.edu/~yipn/421/2024-10-28-ApplsLinProg.pdf#page=3}{\textbf{Production Level}}: The amount of ice cream produced in a given month.<L[2024-10-28-ApplsLinProg] 3><L[2024-10-28-ApplsLinProg] 4><L[2024-10-28-ApplsLinProg] 5>%
\item \href{https://www.math.purdue.edu/~yipn/421/2024-10-28-ApplsLinProg.pdf#page=3}{\textbf{Surplus}}: The excess amount of ice cream remaining at the end of a month after meeting the demand.<L[2024-10-28-ApplsLinProg] 3><L[2024-10-28-ApplsLinProg] 4><L[2024-10-28-ApplsLinProg] 5>%
\item \href{https://www.math.purdue.edu/~yipn/421/2024-10-28-ApplsLinProg.pdf#page=3}{\textbf{Cost of Changing Production}}: The cost associated with adjusting the production level from one month to the next.<L[2024-10-28-ApplsLinProg] 3><L[2024-10-28-ApplsLinProg] 4><L[2024-10-28-ApplsLinProg] 5>%
\item \href{https://www.math.purdue.edu/~yipn/421/2024-10-28-ApplsLinProg.pdf#page=3}{\textbf{Cost of Storage}}: The cost of storing surplus ice cream from one month to the next.<L[2024-10-28-ApplsLinProg] 3><L[2024-10-28-ApplsLinProg] 4><L[2024-10-28-ApplsLinProg] 5>%
\item \href{https://www.math.purdue.edu/~yipn/421/2024-10-28-ApplsLinProg.pdf#page=7}{\textbf{Regular Employees}}: Employees with a fixed hourly wage, used in workforce planning.<L[2024-10-28-ApplsLinProg] 7><L[2024-10-28-ApplsLinProg] 8>%
\item \href{https://www.math.purdue.edu/~yipn/421/2024-10-28-ApplsLinProg.pdf#page=7}{\textbf{Temporary Employees}}: Employees hired on a temporary basis, with a higher hourly wage than regular employees.<L[2024-10-28-ApplsLinProg] 7><L[2024-10-28-ApplsLinProg] 8>%
\item \href{https://www.math.purdue.edu/~yipn/421/2024-10-28-ApplsLinProg.pdf#page=9}{\textbf{Portfolio Selection}}: The problem of choosing the optimal mix of investments to maximize return while managing risk.<L[2024-10-28-ApplsLinProg] 9><L[2024-10-28-ApplsLinProg] 10><L[2024-10-28-ApplsLinProg] 11><L[2024-10-28-ApplsLinProg] 12><L[2024-10-28-ApplsLinProg] 13><L[2024-10-28-ApplsLinProg] 14><L[2024-10-28-ApplsLinProg] 16>%
\item \href{https://www.math.purdue.edu/~yipn/421/2024-10-28-ApplsLinProg.pdf#page=9}{\textbf{Annual Return}}: The return on investment for a given asset over a year.<L[2024-10-28-ApplsLinProg] 9><L[2024-10-28-ApplsLinProg] 10><L[2024-10-28-ApplsLinProg] 11><L[2024-10-28-ApplsLinProg] 12><L[2024-10-28-ApplsLinProg] 13><L[2024-10-28-ApplsLinProg] 14>%
\item \href{https://www.math.purdue.edu/~yipn/421/2024-10-28-ApplsLinProg.pdf#page=9}{\textbf{Proportion of Investment}}: The fraction of a portfolio invested in a particular asset.<L[2024-10-28-ApplsLinProg] 9><L[2024-10-28-ApplsLinProg] 10><L[2024-10-28-ApplsLinProg] 11><L[2024-10-28-ApplsLinProg] 12><L[2024-10-28-ApplsLinProg] 13><L[2024-10-28-ApplsLinProg] 14>%
\item \href{https://www.math.purdue.edu/~yipn/421/2024-10-28-ApplsLinProg.pdf#page=10}{\textbf{Risk}}: The uncertainty or variability associated with the return on investment.<L[2024-10-28-ApplsLinProg] 10><L[2024-10-28-ApplsLinProg] 12><L[2024-10-28-ApplsLinProg] 13><L[2024-10-28-ApplsLinProg] 14>%
\item \href{https://www.math.purdue.edu/~yipn/421/2024-10-28-ApplsLinProg.pdf#page=10}{\textbf{Fluctuation}}: The deviation of an asset's return from its expected return.<L[2024-10-28-ApplsLinProg] 10>%
\item \href{https://www.math.purdue.edu/~yipn/421/2024-10-28-ApplsLinProg.pdf#page=13}{\textbf{Covariance}}: A measure of the joint variability of two random variables, used in portfolio optimization to quantify the relationship between asset returns.<L[2024-10-28-ApplsLinProg] 13><L[2024-10-28-ApplsLinProg] 14>%
\item \href{https://www.math.purdue.edu/~yipn/421/2024-10-28-ApplsLinProg.pdf#page=17}{\textbf{Smallest Enclosing Ball}}: The problem of finding the smallest ball (or hypersphere) that contains a given set of points.<L[2024-10-28-ApplsLinProg] 17><L[2024-10-28-ApplsLinProg] 18>%
\item \href{https://www.math.purdue.edu/~yipn/421/2024-10-28-ApplsLinProg.pdf#page=19}{\textbf{Maximum Weight Matching}}: The problem of finding a matching in a weighted graph that maximizes the total weight of the edges in the matching.<L[2024-10-28-ApplsLinProg] 19><L[2024-10-28-ApplsLinProg] 20><L[2024-10-28-ApplsLinProg] 21>%
\item \href{https://www.math.purdue.edu/~yipn/421/2024-10-28-ApplsLinProg.pdf#page=22}{\textbf{Machine Scheduling}}: The problem of assigning jobs to machines to minimize the total completion time.<L[2024-10-28-ApplsLinProg] 22><L[2024-10-28-ApplsLinProg] 23>%
\item \href{https://www.math.purdue.edu/~yipn/421/2024-10-28-ApplsLinProg.pdf#page=24}{\textbf{Knapsack Problem}}: The problem of selecting a subset of items with weights and values to maximize the total value while not exceeding a weight limit.<L[2024-10-28-ApplsLinProg] 24>%
\item \href{https://www.math.purdue.edu/~yipn/421/2024-10-28-ApplsLinProg.pdf#page=25}{\textbf{Cutting Stock Problem}}: The problem of cutting large rolls of material into smaller rolls of specified widths to minimize waste.<L[2024-10-28-ApplsLinProg] 25><L[2024-10-28-ApplsLinProg] 26><L[2024-10-28-ApplsLinProg] 27><L[2024-10-28-ApplsLinProg] 28><L[2024-10-28-ApplsLinProg] 29>%
\item \href{https://www.math.purdue.edu/~yipn/421/2024-10-28-ApplsLinProg.pdf#page=30}{\textbf{Sparse Solutions of Linear System}}: The problem of finding a solution to a system of linear equations that has a limited number of non-zero elements.<L[2024-10-28-ApplsLinProg] 30><L[2024-10-28-ApplsLinProg] 31><L[2024-10-28-ApplsLinProg] 32><L[2024-10-28-ApplsLinProg] 33><L[2024-10-28-ApplsLinProg] 34><L[2024-10-28-ApplsLinProg] 35><L[2024-10-28-ApplsLinProg] 36><L[2024-10-28-ApplsLinProg] 37>%
\item \href{https://www.math.purdue.edu/~yipn/421/2024-10-28-ApplsLinProg.pdf#page=34}{\textbf{\$l\_1\$ Norm}}: The sum of the absolute values of the elements of a vector, used as an objective function to find sparse solutions.<L[2024-10-28-ApplsLinProg] 34><L[2024-10-28-ApplsLinProg] 35><L[2024-10-28-ApplsLinProg] 36><L[2024-10-28-ApplsLinProg] 37>%
\end{itemize}%
\begin{itemize}[[leftmargin=*]]%
\item \href{https://www.math.purdue.edu/~yipn/421/2024-10-28-Farkas.pdf#page=1}{\textbf{Farkas Lemma}}: A fundamental theorem in linear programming that provides conditions for the existence or non-existence of solutions to systems of linear inequalities.<L[2024-10-28-Farkas] 1><L[2024-10-28-Farkas] 2><L[2024-10-28-Farkas] 3><L[2024-10-28-Farkas] 4><L[2024-10-28-Farkas] 5><L[2024-10-28-Farkas] 6><L[2024-10-28-Farkas] 7><L[2024-10-28-Farkas] 8><L[2024-10-28-Farkas] 9><L[2024-10-28-Farkas] 10><L[2024-10-28-Farkas] 11>%
\item \href{https://www.math.purdue.edu/~yipn/421/2024-10-28-Farkas.pdf#page=3}{\textbf{Fredholm Alternatives}}: A statement of the conditions under which a system of linear equations $Ax = b$ has a solution, and the conditions under which a related homogeneous system has a non-trivial solution.<L[2024-10-28-Farkas] 3>%
\item \href{https://www.math.purdue.edu/~yipn/421/2024-10-28-Farkas.pdf#page=1}{\textbf{Primal Linear Program}}: A linear programming problem in its standard form, typically aiming to maximize an objective function subject to constraints.<L[2024-10-28-Farkas] 1>%
\item \href{https://www.math.purdue.edu/~yipn/421/2024-10-28-Farkas.pdf#page=1}{\textbf{Dual Linear Program}}: A linear programming problem derived from the primal problem, providing a lower bound (for maximization problems) on the optimal value of the primal problem.<L[2024-10-28-Farkas] 1>%
\end{itemize}%
\begin{itemize}[[leftmargin=*]]%
\item \href{https://www.math.purdue.edu/~yipn/421/2024-10-31-IntLinProg.pdf#page=1}{\textbf{Integer Programming}}: A mathematical optimization problem where some or all of the variables are restricted to integer values.<L[2024-10-31-IntLinProg] 1><L[2024-10-31-IntLinProg] 2><L[2024-10-31-IntLinProg] 3><L[2024-10-31-IntLinProg] 4><L[2024-10-31-IntLinProg] 5><L[2024-10-31-IntLinProg] 6><L[2024-10-31-IntLinProg] 7><L[2024-10-31-IntLinProg] 8><L[2024-10-31-IntLinProg] 9><L[2024-10-31-IntLinProg] 10><L[2024-10-31-IntLinProg] 15><L[2024-10-31-IntLinProg] 16><L[2024-10-31-IntLinProg] 17><L[2024-10-31-IntLinProg] 18><L[2024-10-31-IntLinProg] 19><L[2024-10-31-IntLinProg] 20><L[2024-10-31-IntLinProg] 21><L[2024-10-31-IntLinProg] 22><L[2024-10-31-IntLinProg] 23>%
\item \href{https://www.math.purdue.edu/~yipn/421/2024-10-31-IntLinProg.pdf#page=2}{\textbf{Relaxed Problem}}: The linear programming problem obtained by removing the integer constraints from an integer programming problem.<L[2024-10-31-IntLinProg] 2><L[2024-10-31-IntLinProg] 4><L[2024-10-31-IntLinProg] 5><L[2024-10-31-IntLinProg] 6><L[2024-10-31-IntLinProg] 7><L[2024-10-31-IntLinProg] 8><L[2024-10-31-IntLinProg] 9><L[2024-10-31-IntLinProg] 10><L[2024-10-31-IntLinProg] 16><L[2024-10-31-IntLinProg] 17><L[2024-10-31-IntLinProg] 18><L[2024-10-31-IntLinProg] 19><L[2024-10-31-IntLinProg] 20><L[2024-10-31-IntLinProg] 21><L[2024-10-31-IntLinProg] 22>%
\item \href{https://www.math.purdue.edu/~yipn/421/2024-10-31-IntLinProg.pdf#page=5}{\textbf{Branch and Bound}}: A method for solving integer programming problems by recursively partitioning the feasible region and bounding the objective function value.<L[2024-10-31-IntLinProg] 5><L[2024-10-31-IntLinProg] 6><L[2024-10-31-IntLinProg] 7><L[2024-10-31-IntLinProg] 8><L[2024-10-31-IntLinProg] 9><L[2024-10-31-IntLinProg] 10><L[2024-10-31-IntLinProg] 11><L[2024-10-31-IntLinProg] 12><L[2024-10-31-IntLinProg] 13><L[2024-10-31-IntLinProg] 14>%
\item \href{https://www.math.purdue.edu/~yipn/421/2024-10-31-IntLinProg.pdf#page=15}{\textbf{Gomory Cuts}}: Additional constraints added to a linear programming relaxation of an integer program to cut off non-integer solutions.<L[2024-10-31-IntLinProg] 15><L[2024-10-31-IntLinProg] 16><L[2024-10-31-IntLinProg] 17><L[2024-10-31-IntLinProg] 18><L[2024-10-31-IntLinProg] 19><L[2024-10-31-IntLinProg] 20><L[2024-10-31-IntLinProg] 21><L[2024-10-31-IntLinProg] 22><L[2024-10-31-IntLinProg] 23>%
\end{itemize}%
\begin{itemize}[[leftmargin=*]]%
\item \href{https://www.math.purdue.edu/~yipn/421/2024-11-05-NetFlows.pdf#page=1}{\textbf{Network}}: A graph consisting of a set of nodes and a set of directed arcs connecting the nodes.<L[2024-11-05-NetFlows] 1>%
\item \href{https://www.math.purdue.edu/~yipn/421/2024-11-05-NetFlows.pdf#page=1}{\textbf{Arc}}: A directed edge connecting two nodes in a network, often associated with a capacity or cost.<L[2024-11-05-NetFlows] 1>%
\item \href{https://www.math.purdue.edu/~yipn/421/2024-11-05-NetFlows.pdf#page=1}{\textbf{Node}}: A point or vertex in a network representing a location or entity.<L[2024-11-05-NetFlows] 1>%
\item \href{https://www.math.purdue.edu/~yipn/421/2024-11-05-NetFlows.pdf#page=2}{\textbf{Supply}}: The net flow at a node, represented by a positive value of $b_i$.<L[2024-11-05-NetFlows] 2>%
\item \href{https://www.math.purdue.edu/~yipn/421/2024-11-05-NetFlows.pdf#page=2}{\textbf{Demand}}: The net flow at a node, represented by a negative value of $b_i$.<L[2024-11-05-NetFlows] 2>%
\item \href{https://www.math.purdue.edu/~yipn/421/2024-11-05-NetFlows.pdf#page=2}{\textbf{\$b\_i\$}}: The net flow (supply or demand) at node $i$ in a network flow problem. $\sum_{i \in N} b_i = 0$.<L[2024-11-05-NetFlows] 2>%
\item \href{https://www.math.purdue.edu/~yipn/421/2024-11-05-NetFlows.pdf#page=2}{\textbf{\$X\_\{ij\}\$}}: The amount of flow transported from node $i$ to node $j$ in a network.<L[2024-11-05-NetFlows] 2>%
\item \href{https://www.math.purdue.edu/~yipn/421/2024-11-05-NetFlows.pdf#page=2}{\textbf{\$C\_\{ij\}\$}}: The cost of transporting one unit of flow from node $i$ to node $j$ in a network.<L[2024-11-05-NetFlows] 2>%
\item \href{https://www.math.purdue.edu/~yipn/421/2024-11-05-NetFlows.pdf#page=3}{\textbf{Balanced Equation}}: An equation stating that the sum of flows entering a node plus the supply at that node equals the sum of flows leaving the node.<L[2024-11-05-NetFlows] 3>%
\item \href{https://www.math.purdue.edu/~yipn/421/2024-11-05-NetFlows.pdf#page=4}{\textbf{Incidence Matrix (A)}}: A matrix representing the connections between nodes and arcs in a network. Each row represents a node, each column an arc.<L[2024-11-05-NetFlows] 4>%
\item \href{https://www.math.purdue.edu/~yipn/421/2024-11-05-NetFlows.pdf#page=4}{\textbf{Cost Vector (c)}}: A vector containing the cost of transporting one unit of flow along each arc.<L[2024-11-05-NetFlows] 4>%
\item \href{https://www.math.purdue.edu/~yipn/421/2024-11-05-NetFlows.pdf#page=4}{\textbf{Supply/Demand Vector (b)}}: A vector representing the net flow (supply or demand) at each node.<L[2024-11-05-NetFlows] 4>%
\item \href{https://www.math.purdue.edu/~yipn/421/2024-11-05-NetFlows.pdf#page=5}{\textbf{\$A\_\{ij\}(k,l)\$}}: An element of the incidence matrix defining the value based on the relationship between row index ($i$) and column indices ($k$ and $l$).<L[2024-11-05-NetFlows] 5>%
\item \href{https://www.math.purdue.edu/~yipn/421/2024-11-05-NetFlows.pdf#page=6}{\textbf{Rank(A)}}: The rank of the incidence matrix, which is equal to $m-1$ for a connected network, where $m$ is the number of arcs.<L[2024-11-05-NetFlows] 6>%
\item \href{https://www.math.purdue.edu/~yipn/421/2024-11-05-NetFlows.pdf#page=7}{\textbf{Primal Problem (P)}}: The original minimization problem in linear programming, aiming to minimize the total cost subject to flow conservation and non-negativity constraints.<L[2024-11-05-NetFlows] 7>%
\item \href{https://www.math.purdue.edu/~yipn/421/2024-11-05-NetFlows.pdf#page=7}{\textbf{Dual Problem (D)}}: The maximization problem in linear programming, related to the primal problem through duality.<L[2024-11-05-NetFlows] 7>%
\item \href{https://www.math.purdue.edu/~yipn/421/2024-11-05-NetFlows.pdf#page=10}{\textbf{Path}}: A sequence of nodes in a network where each consecutive pair is connected by an arc.<L[2024-11-05-NetFlows] 10>%
\item \href{https://www.math.purdue.edu/~yipn/421/2024-11-05-NetFlows.pdf#page=10}{\textbf{Cycle}}: A path that starts and ends at the same node.<L[2024-11-05-NetFlows] 10>%
\item \href{https://www.math.purdue.edu/~yipn/421/2024-11-05-NetFlows.pdf#page=10}{\textbf{Tree}}: A connected graph with no cycles.<L[2024-11-05-NetFlows] 10>%
\item \href{https://www.math.purdue.edu/~yipn/421/2024-11-05-NetFlows.pdf#page=10}{\textbf{Spanning Tree}}: A tree that includes all nodes in a network.<L[2024-11-05-NetFlows] 10>%
\item \href{https://www.math.purdue.edu/~yipn/421/2024-11-05-NetFlows.pdf#page=11}{\textbf{Balanced Flow}}: A flow satisfying the flow conservation constraints ($AX = -b$).<L[2024-11-05-NetFlows] 11>%
\item \href{https://www.math.purdue.edu/~yipn/421/2024-11-05-NetFlows.pdf#page=11}{\textbf{Feasible Flow}}: A flow satisfying both flow conservation and non-negativity constraints ($X \ge 0$).<L[2024-11-05-NetFlows] 11>%
\item \href{https://www.math.purdue.edu/~yipn/421/2024-11-05-NetFlows.pdf#page=11}{\textbf{Tree Solution}}: A flow solution where the flow along arcs in the spanning tree is zero ($X(i,j) = 0$ if $(i,j) \in \mathcal{T}$).<L[2024-11-05-NetFlows] 11>%
\item \href{https://www.math.purdue.edu/~yipn/421/2024-11-05-NetFlows.pdf#page=16}{\textbf{Root Node}}: A node selected to simplify the representation of the flow equations by deleting its corresponding row from the incidence matrix.<L[2024-11-05-NetFlows] 16>%
\item \href{https://www.math.purdue.edu/~yipn/421/2024-11-05-NetFlows.pdf#page=15}{\textbf{Basis}}: A square submatrix of the modified incidence matrix ($\tilde{A}$) corresponding to a spanning tree.<L[2024-11-05-NetFlows] 15>%
\item \href{https://www.math.purdue.edu/~yipn/421/2024-11-05-NetFlows.pdf#page=18}{\textbf{Complementary Slackness}}: A condition stating that $x_{ij}z_{ij} = 0$ for all arcs $(i,j)$ in a network flow problem.<L[2024-11-05-NetFlows] 18>%
\item \href{https://www.math.purdue.edu/~yipn/421/2024-11-05-NetFlows.pdf#page=18}{\textbf{Dual Flow}}: The flow in the dual problem, represented by node potentials ($y_i$) and dual slack variables ($z_{ij}$).<L[2024-11-05-NetFlows] 18>%
\end{itemize}%
\begin{itemize}[[leftmargin=*]]%
\item \href{https://www.math.purdue.edu/~yipn/421/2024-11-07-NetSimplex.pdf#page=1}{\textbf{Network}}: A graph representing a system of nodes and edges, often used to model flow problems. <L[2024-11-07-NetSimplex] 1><L[2024-11-07-NetSimplex] 2><L[2024-11-07-NetSimplex] 3><L[2024-11-07-NetSimplex] 4>%
\item \href{https://www.math.purdue.edu/~yipn/421/2024-11-07-NetSimplex.pdf#page=1}{\textbf{Spanning Tree}}: A subgraph of a network that includes all nodes and is connected and acyclic (contains no cycles). <L[2024-11-07-NetSimplex] 1><L[2024-11-07-NetSimplex] 2><L[2024-11-07-NetSimplex] 3><L[2024-11-07-NetSimplex] 4>%
\item \href{https://www.math.purdue.edu/~yipn/421/2024-11-07-NetSimplex.pdf#page=2}{\textbf{Primal Flow}}: A feasible solution to the primal network flow problem, satisfying flow conservation constraints. <L[2024-11-07-NetSimplex] 2>%
\item \href{https://www.math.purdue.edu/~yipn/421/2024-11-07-NetSimplex.pdf#page=3}{\textbf{Dual Flow}}: A feasible solution to the dual network flow problem, satisfying dual constraints. <L[2024-11-07-NetSimplex] 3>%
\item \href{https://www.math.purdue.edu/~yipn/421/2024-11-07-NetSimplex.pdf#page=6}{\textbf{Entering Arc}}: An edge added to the spanning tree during a pivot step in the network simplex method. <L[2024-11-07-NetSimplex] 6><L[2024-11-07-NetSimplex] 7><L[2024-11-07-NetSimplex] 8>%
\item \href{https://www.math.purdue.edu/~yipn/421/2024-11-07-NetSimplex.pdf#page=6}{\textbf{Leaving Arc}}: An edge removed from the spanning tree during a pivot step in the network simplex method. <L[2024-11-07-NetSimplex] 6><L[2024-11-07-NetSimplex] 7><L[2024-11-07-NetSimplex] 8>%
\item \href{https://www.math.purdue.edu/~yipn/421/2024-11-07-NetSimplex.pdf#page=3}{\textbf{Dual Slack}}: The difference between the cost of an arc and the difference in dual variables between its endpoints. <L[2024-11-07-NetSimplex] 3><L[2024-11-07-NetSimplex] 35><L[2024-11-07-NetSimplex] 36><L[2024-11-07-NetSimplex] 37><L[2024-11-07-NetSimplex] 38><L[2024-11-07-NetSimplex] 39><L[2024-11-07-NetSimplex] 40><L[2024-11-07-NetSimplex] 41><L[2024-11-07-NetSimplex] 42>%
\item \href{https://www.math.purdue.edu/~yipn/421/2024-11-07-NetSimplex.pdf#page=34}{\textbf{Subtrees}}: The two disconnected trees resulting from removing a leaving arc from the spanning tree. <L[2024-11-07-NetSimplex] 34><L[2024-11-07-NetSimplex] 35><L[2024-11-07-NetSimplex] 36><L[2024-11-07-NetSimplex] 37><L[2024-11-07-NetSimplex] 38><L[2024-11-07-NetSimplex] 39><L[2024-11-07-NetSimplex] 40><L[2024-11-07-NetSimplex] 41><L[2024-11-07-NetSimplex] 42>%
\end{itemize}%
\begin{itemize}[[leftmargin=*]]%
\item \href{https://www.math.purdue.edu/~yipn/421/2024-11-14-NetAppls.pdf#page=1}{\textbf{Transportation Problem}}: A linear programming problem where the goal is to minimize the cost of transporting goods from multiple sources to multiple destinations, subject to supply and demand constraints. <L[2024-11-14-NetAppls] 1><L[2024-11-14-NetAppls] 2><L[2024-11-14-NetAppls] 3><L[2024-11-14-NetAppls] 4><L[2024-11-14-NetAppls] 5><L[2024-11-14-NetAppls] 6>%
\item \href{https://www.math.purdue.edu/~yipn/421/2024-11-14-NetAppls.pdf#page=3}{\textbf{Dummy Node}}: A node added to a network flow model to handle excess supply or unmet demand, converting inequality constraints into equalities. <L[2024-11-14-NetAppls] 3><L[2024-11-14-NetAppls] 4><L[2024-11-14-NetAppls] 8>%
\item \href{https://www.math.purdue.edu/~yipn/421/2024-11-14-NetAppls.pdf#page=9}{\textbf{Upper Bounded Transshipment Problem}}: A network flow problem where the flow along each arc is constrained by an upper bound in addition to supply and demand constraints. <L[2024-11-14-NetAppls] 9><L[2024-11-14-NetAppls] 10><L[2024-11-14-NetAppls] 11><L[2024-11-14-NetAppls] 12>%
\item \href{https://www.math.purdue.edu/~yipn/421/2024-11-14-NetAppls.pdf#page=13}{\textbf{Shortest Path Problem}}: A graph problem where the goal is to find the path with the minimum total weight or cost between two specified nodes in a weighted graph. <L[2024-11-14-NetAppls] 13><L[2024-11-14-NetAppls] 14><L[2024-11-14-NetAppls] 15><L[2024-11-14-NetAppls] 16><L[2024-11-14-NetAppls] 17><L[2024-11-14-NetAppls] 18><L[2024-11-14-NetAppls] 19><L[2024-11-14-NetAppls] 20><L[2024-11-14-NetAppls] 21>%
\item \href{https://www.math.purdue.edu/~yipn/421/2024-11-14-NetAppls.pdf#page=16}{\textbf{Bellman's Equation}}: A recursive equation used to solve the shortest path problem by iteratively computing the shortest distance from each node to a destination node. <L[2024-11-14-NetAppls] 16><L[2024-11-14-NetAppls] 17><L[2024-11-14-NetAppls] 18><L[2024-11-14-NetAppls] 19>%
\item \href{https://www.math.purdue.edu/~yipn/421/2024-11-14-NetAppls.pdf#page=20}{\textbf{Dijkstra's Algorithm}}: A greedy algorithm that finds the shortest paths from a single source node to all other nodes in a graph with non-negative edge weights. <L[2024-11-14-NetAppls] 20><L[2024-11-14-NetAppls] 21>%
\end{itemize}%
\begin{itemize}[[leftmargin=*]]%
\item \href{https://www.math.purdue.edu/~yipn/421/2024-11-19-MaxFlowMinCut.pdf#page=1}{\textbf{Max Flow}}: The maximum amount of flow that can be sent from a source node to a sink node in a network, subject to capacity constraints on the edges. <L[2024-11-19-MaxFlowMinCut] 1><L[2024-11-19-MaxFlowMinCut] 2><L[2024-11-19-MaxFlowMinCut] 3><L[2024-11-19-MaxFlowMinCut] 4><L[2024-11-19-MaxFlowMinCut] 5><L[2024-11-19-MaxFlowMinCut] 6><L[2024-11-19-MaxFlowMinCut] 7><L[2024-11-19-MaxFlowMinCut] 8><L[2024-11-19-MaxFlowMinCut] 9><L[2024-11-19-MaxFlowMinCut] 10><L[2024-11-19-MaxFlowMinCut] 11><L[2024-11-19-MaxFlowMinCut] 12><L[2024-11-19-MaxFlowMinCut] 13><L[2024-11-19-MaxFlowMinCut] 14><L[2024-11-19-MaxFlowMinCut] 15><L[2024-11-19-MaxFlowMinCut] 16>%
\item \href{https://www.math.purdue.edu/~yipn/421/2024-11-19-MaxFlowMinCut.pdf#page=2}{\textbf{Min Cut}}: A partition of the nodes in a network into two sets, one containing the source node and the other containing the sink node. The capacity of the min cut is the sum of the capacities of the edges crossing from the source set to the sink set. <L[2024-11-19-MaxFlowMinCut] 2><L[2024-11-19-MaxFlowMinCut] 3><L[2024-11-19-MaxFlowMinCut] 4><L[2024-11-19-MaxFlowMinCut] 8><L[2024-11-19-MaxFlowMinCut] 9><L[2024-11-19-MaxFlowMinCut] 14><L[2024-11-19-MaxFlowMinCut] 15><L[2024-11-19-MaxFlowMinCut] 16>%
\item \href{https://www.math.purdue.edu/~yipn/421/2024-11-19-MaxFlowMinCut.pdf#page=2}{\textbf{Cut Capacity}}: The sum of the capacities of all edges that cross a cut from the source set to the sink set. <L[2024-11-19-MaxFlowMinCut] 2><L[2024-11-19-MaxFlowMinCut] 3><L[2024-11-19-MaxFlowMinCut] 4><L[2024-11-19-MaxFlowMinCut] 8><L[2024-11-19-MaxFlowMinCut] 9>%
\item \href{https://www.math.purdue.edu/~yipn/421/2024-11-19-MaxFlowMinCut.pdf#page=4}{\textbf{Max{-}Flow Min{-}Cut Theorem}}: A fundamental theorem in network flow theory stating that the maximum flow from a source to a sink in a network is equal to the minimum capacity of any cut separating the source and sink. <L[2024-11-19-MaxFlowMinCut] 4><L[2024-11-19-MaxFlowMinCut] 8><L[2024-11-19-MaxFlowMinCut] 9><L[2024-11-19-MaxFlowMinCut] 16>%
\item \href{https://www.math.purdue.edu/~yipn/421/2024-11-19-MaxFlowMinCut.pdf#page=10}{\textbf{Augmenting Path}}: A path in a network from the source to the sink along which additional flow can be sent without violating capacity constraints. <L[2024-11-19-MaxFlowMinCut] 10><L[2024-11-19-MaxFlowMinCut] 11><L[2024-11-19-MaxFlowMinCut] 12><L[2024-11-19-MaxFlowMinCut] 13>%
\item \href{https://www.math.purdue.edu/~yipn/421/2024-11-19-MaxFlowMinCut.pdf#page=10}{\textbf{Ford{-}Fulkerson Algorithm}}: An algorithm for finding the maximum flow in a network by iteratively finding augmenting paths and increasing the flow along them. <L[2024-11-19-MaxFlowMinCut] 10><L[2024-11-19-MaxFlowMinCut] 11><L[2024-11-19-MaxFlowMinCut] 12><L[2024-11-19-MaxFlowMinCut] 13>%
\item \href{https://www.math.purdue.edu/~yipn/421/2024-11-19-MaxFlowMinCut.pdf#page=1}{\textbf{Network Flow Problem}}: A problem of finding the maximum flow that can be sent through a network from a source node to a sink node, subject to capacity constraints on the edges. <L[2024-11-19-MaxFlowMinCut] 1><L[2024-11-19-MaxFlowMinCut] 2><L[2024-11-19-MaxFlowMinCut] 3><L[2024-11-19-MaxFlowMinCut] 4><L[2024-11-19-MaxFlowMinCut] 5><L[2024-11-19-MaxFlowMinCut] 6><L[2024-11-19-MaxFlowMinCut] 7><L[2024-11-19-MaxFlowMinCut] 8><L[2024-11-19-MaxFlowMinCut] 9><L[2024-11-19-MaxFlowMinCut] 10><L[2024-11-19-MaxFlowMinCut] 11><L[2024-11-19-MaxFlowMinCut] 12><L[2024-11-19-MaxFlowMinCut] 13><L[2024-11-19-MaxFlowMinCut] 14><L[2024-11-19-MaxFlowMinCut] 15><L[2024-11-19-MaxFlowMinCut] 16>%
\item \href{https://www.math.purdue.edu/~yipn/421/2024-11-19-MaxFlowMinCut.pdf#page=1}{\textbf{Source Node}}: The node in a network from which flow originates. <L[2024-11-19-MaxFlowMinCut] 1><L[2024-11-19-MaxFlowMinCut] 2><L[2024-11-19-MaxFlowMinCut] 3><L[2024-11-19-MaxFlowMinCut] 4><L[2024-11-19-MaxFlowMinCut] 8><L[2024-11-19-MaxFlowMinCut] 9><L[2024-11-19-MaxFlowMinCut] 14><L[2024-11-19-MaxFlowMinCut] 15>%
\item \href{https://www.math.purdue.edu/~yipn/421/2024-11-19-MaxFlowMinCut.pdf#page=1}{\textbf{Sink Node}}: The node in a network to which flow is destined. <L[2024-11-19-MaxFlowMinCut] 1><L[2024-11-19-MaxFlowMinCut] 2><L[2024-11-19-MaxFlowMinCut] 3><L[2024-11-19-MaxFlowMinCut] 4><L[2024-11-19-MaxFlowMinCut] 8><L[2024-11-19-MaxFlowMinCut] 9><L[2024-11-19-MaxFlowMinCut] 14><L[2024-11-19-MaxFlowMinCut] 15>%
\item \href{https://www.math.purdue.edu/~yipn/421/2024-11-19-MaxFlowMinCut.pdf#page=1}{\textbf{Edge Capacity}}: The maximum amount of flow that can pass through an edge in a network. <L[2024-11-19-MaxFlowMinCut] 1><L[2024-11-19-MaxFlowMinCut] 2><L[2024-11-19-MaxFlowMinCut] 3><L[2024-11-19-MaxFlowMinCut] 4><L[2024-11-19-MaxFlowMinCut] 5><L[2024-11-19-MaxFlowMinCut] 6><L[2024-11-19-MaxFlowMinCut] 7><L[2024-11-19-MaxFlowMinCut] 8><L[2024-11-19-MaxFlowMinCut] 9><L[2024-11-19-MaxFlowMinCut] 14><L[2024-11-19-MaxFlowMinCut] 15><L[2024-11-19-MaxFlowMinCut] 16>%
\end{itemize}%
\begin{itemize}[[leftmargin=*]]%
\item \href{https://www.math.purdue.edu/~yipn/421/2024-11-21-EconNetSimplex.pdf#page=1}{\textbf{Network Simplex Method}}: A specialized simplex algorithm for solving linear programming problems with network structures, such as transportation or assignment problems. <L[2024-11-21-EconNetSimplex] 1><L[2024-11-21-EconNetSimplex] 2><L[2024-11-21-EconNetSimplex] 3>%
\item \href{https://www.math.purdue.edu/~yipn/421/2024-11-21-EconNetSimplex.pdf#page=3}{\textbf{Spanning Tree}}: A subset of edges in a network graph that connects all nodes without forming any cycles. <L[2024-11-21-EconNetSimplex] 3>%
\item \href{https://www.math.purdue.edu/~yipn/421/2024-11-21-EconNetSimplex.pdf#page=3}{\textbf{Feasible Solution}}: A solution to a linear programming problem that satisfies all constraints. <L[2024-11-21-EconNetSimplex] 3><L[2024-11-21-EconNetSimplex] 4>%
\item \href{https://www.math.purdue.edu/~yipn/421/2024-11-21-EconNetSimplex.pdf#page=4}{\textbf{Dual Prices (or Shadow Prices)}}: In the context of network simplex, these represent the "fair" market prices at each node, reflecting the cost of shipping and maintaining profitability. <L[2024-11-21-EconNetSimplex] 4>%
\item \href{https://www.math.purdue.edu/~yipn/421/2024-11-21-EconNetSimplex.pdf#page=5}{\textbf{Reduced Cost}}: The difference between the cost of shipping along an arc and the difference in dual prices between the nodes connected by that arc ($y_i + c_{ij} - y_j$). A negative reduced cost indicates a potential improvement in the solution. <L[2024-11-21-EconNetSimplex] 5>%
\end{itemize}%
\begin{itemize}[[leftmargin=*]]%
\item \href{https://www.math.purdue.edu/~yipn/421/2024-12-03-Duality.pdf#page=1}{\textbf{Duality}}: In linear programming, duality refers to the relationship between a primal problem and its corresponding dual problem. The primal problem seeks to maximize (or minimize) an objective function subject to constraints, while the dual problem seeks to minimize (or maximize) a different objective function subject to different constraints derived from the primal problem. The optimal solutions of both problems are related, and the dual problem provides valuable insights into the primal problem. <L[2024-12-03-Duality] 1>%
\item \href{https://www.math.purdue.edu/~yipn/421/2024-12-03-Duality.pdf#page=2}{\textbf{Complementary Slackness}}: At the optimal solution of a linear programming problem and its dual, complementary slackness states that for each constraint in the primal problem, either the constraint is binding (the slack variable is zero), or the corresponding dual variable is zero (and vice versa). This relationship holds for both the primal and dual problems. <L[2024-12-03-Duality] 2>%
\item \href{https://www.math.purdue.edu/~yipn/421/2024-12-03-Duality.pdf#page=3}{\textbf{Graphical Method}}: A method for solving linear programming problems with two variables by plotting the constraints as lines on a graph. The feasible region is the area satisfying all constraints. The optimal solution is found at a corner point of the feasible region that maximizes or minimizes the objective function. <L[2024-12-03-Duality] 3><L[2024-12-03-Duality] 4>%
\item \href{https://www.math.purdue.edu/~yipn/421/2024-12-03-Duality.pdf#page=11}{\textbf{Sensitivity Analysis}}: The study of how changes in the parameters of a linear programming problem (objective function coefficients or constraint values) affect the optimal solution. It helps determine the range of parameter values for which the current optimal solution remains optimal. <L[2024-12-03-Duality] 11><L[2024-12-03-Duality] 12><L[2024-12-03-Duality] 13><L[2024-12-03-Duality] 14><L[2024-12-03-Duality] 15><L[2024-12-03-Duality] 20><L[2024-12-03-Duality] 21>%
\item \href{https://www.math.purdue.edu/~yipn/421/2024-12-03-Duality.pdf#page=15}{\textbf{Shadow Prices}}: In linear programming, shadow prices (also known as dual values) represent the marginal value of increasing a resource constraint by one unit. They are the optimal values of the dual variables and indicate how much the optimal objective function value would improve if the corresponding constraint were relaxed. <L[2024-12-03-Duality] 15><L[2024-12-03-Duality] 20><L[2024-12-03-Duality] 21>%
\end{itemize}%
\newpage

%
\section{Problem Types}%
\label{sec:ProblemTypes}%
\begin{itemize}[[leftmargin=*]]%
\item \href{https://www.math.purdue.edu/~yipn/421/2024-08-27-ExSimplex.pdf#page=32}{\textbf{Finding an Initial Feasible Solution}}: The problem of finding a starting feasible solution for the Simplex method, particularly when the origin is not feasible. An auxiliary problem is introduced to find a feasible solution. <L[2024-08-27-ExSimplex] 32><L[2024-08-27-ExSimplex] 33><L[2024-08-27-ExSimplex] 34><L[2024-08-27-ExSimplex] 35><L[2024-08-27-ExSimplex] 36><L[2024-08-27-ExSimplex] 37><L[2024-08-27-ExSimplex] 38><L[2024-08-27-ExSimplex] 39><L[2024-08-27-ExSimplex] 40><L[2024-08-27-ExSimplex] 41><L[2024-08-27-ExSimplex] 42>%
\item \href{https://www.math.purdue.edu/~yipn/421/2024-08-27-ExSimplex.pdf#page=2}{\textbf{Graphical Representation of Linear Inequalities and Feasible Regions}}: The problem of visually representing linear inequality constraints in two dimensions and identifying the feasible region (the set of points satisfying all constraints). <L[2024-08-27-ExSimplex] 2><L[2024-08-27-ExSimplex] 3><L[2024-08-27-ExSimplex] 4><L[2024-08-27-ExSimplex] 5><L[2024-08-27-ExSimplex] 16>%
\item \href{https://www.math.purdue.edu/~yipn/421/2024-08-27-ExSimplex.pdf#page=4}{\textbf{Solving Linear Programming Problems Graphically}}: The problem of finding the optimal solution to a linear programming problem by graphically identifying the vertices of the feasible region and evaluating the objective function at each vertex. <L[2024-08-27-ExSimplex] 4><L[2024-08-27-ExSimplex] 5><L[2024-08-27-ExSimplex] 16>%
\item \href{https://www.math.purdue.edu/~yipn/421/2024-08-27-ExSimplex.pdf#page=6}{\textbf{Applying the Simplex Method}}: The problem of using the Simplex algorithm to iteratively find the optimal solution to a linear programming problem by moving from one vertex of the feasible region to another, improving the objective function value at each step. <L[2024-08-27-ExSimplex] 6><L[2024-08-27-ExSimplex] 7><L[2024-08-27-ExSimplex] 8><L[2024-08-27-ExSimplex] 9><L[2024-08-27-ExSimplex] 10><L[2024-08-27-ExSimplex] 11><L[2024-08-27-ExSimplex] 12><L[2024-08-27-ExSimplex] 13><L[2024-08-27-ExSimplex] 14><L[2024-08-27-ExSimplex] 15><L[2024-08-27-ExSimplex] 17><L[2024-08-27-ExSimplex] 18><L[2024-08-27-ExSimplex] 19><L[2024-08-27-ExSimplex] 20><L[2024-08-27-ExSimplex] 21><L[2024-08-27-ExSimplex] 22><L[2024-08-27-ExSimplex] 23><L[2024-08-27-ExSimplex] 24><L[2024-08-27-ExSimplex] 25>%
\item \href{https://www.math.purdue.edu/~yipn/421/2024-08-27-ExSimplex.pdf#page=26}{\textbf{Working with Dictionaries of Variables in the Simplex Method}}: The problem of representing and manipulating the linear programming problem using a dictionary of variables, where basic variables are expressed in terms of non-basic variables. This is crucial for the pivot operations in the Simplex method. <L[2024-08-27-ExSimplex] 26><L[2024-08-27-ExSimplex] 27><L[2024-08-27-ExSimplex] 28><L[2024-08-27-ExSimplex] 29>%
\item \href{https://www.math.purdue.edu/~yipn/421/2024-08-27-ExSimplex.pdf#page=30}{\textbf{Handling Infeasible Origins}}: The problem of dealing with linear programming problems where the origin (0,0) is not a feasible solution. This requires the use of an auxiliary problem to find a feasible starting point. <L[2024-08-27-ExSimplex] 30><L[2024-08-27-ExSimplex] 31><L[2024-08-27-ExSimplex] 32><L[2024-08-27-ExSimplex] 33><L[2024-08-27-ExSimplex] 34><L[2024-08-27-ExSimplex] 35><L[2024-08-27-ExSimplex] 36><L[2024-08-27-ExSimplex] 37><L[2024-08-27-ExSimplex] 38><L[2024-08-27-ExSimplex] 39><L[2024-08-27-ExSimplex] 40><L[2024-08-27-ExSimplex] 41><L[2024-08-27-ExSimplex] 42><L[2024-08-27-ExSimplex] 43>%
\end{itemize}%
\begin{itemize}[[leftmargin=*]]%
\item \href{https://www.math.purdue.edu/~yipn/421/2024-09-05-SimplexEfficiency.pdf#page=2}{\textbf{Worst{-}Case Analysis of the Simplex Method}}: Analyzing the performance of the Simplex method in scenarios designed to maximize the number of iterations required to reach an optimal solution. The Klee-Minty example is a prime example of this type of analysis. <L[2024-09-05-SimplexEfficiency] 2><L[2024-09-05-SimplexEfficiency] 3><L[2024-09-05-SimplexEfficiency] 4><L[2024-09-05-SimplexEfficiency] 5>%
\item \href{https://www.math.purdue.edu/~yipn/421/2024-09-05-SimplexEfficiency.pdf#page=1}{\textbf{Average{-}Case Analysis of the Simplex Method}}: Analyzing the average number of iterations required by the Simplex method across a range of problem instances. This analysis often involves empirical studies and simulations. <L[2024-09-05-SimplexEfficiency] 1><L[2024-09-05-SimplexEfficiency] 7>%
\item \href{https://www.math.purdue.edu/~yipn/421/2024-09-05-SimplexEfficiency.pdf#page=3}{\textbf{Impact of Pivot Rules on Simplex Method Efficiency}}: Investigating how different rules for selecting the entering and leaving variables in the Simplex method affect the number of iterations required to find a solution. The largest-coefficient rule and the largest-increase rule are compared in this context. <L[2024-09-05-SimplexEfficiency] 3><L[2024-09-05-SimplexEfficiency] 6><L[2024-09-05-SimplexEfficiency] 7><L[2024-09-05-SimplexEfficiency] 8>%
\item \href{https://www.math.purdue.edu/~yipn/421/2024-09-05-SimplexEfficiency.pdf#page=5}{\textbf{Sensitivity of Simplex Method to Variable Scaling}}: Examining how scaling the variables in a linear program affects the performance of the Simplex method, particularly with respect to the choice of pivot rule. <L[2024-09-05-SimplexEfficiency] 5>%
\item \href{https://www.math.purdue.edu/~yipn/421/2024-09-05-SimplexEfficiency.pdf#page=9}{\textbf{Empirical Studies of Simplex Method Performance}}: Analyzing the performance of the Simplex method through empirical studies and simulations, often involving large-scale problem instances. The results are often presented graphically or in tabular form. <L[2024-09-05-SimplexEfficiency] 9>%
\item \href{https://www.math.purdue.edu/~yipn/421/2024-09-05-SimplexEfficiency.pdf#page=10}{\textbf{Theoretical Bounds on Simplex Method Iterations}}: Investigating theoretical upper bounds on the number of iterations required by the Simplex method, often involving randomized pivot rules. The focus is on whether polynomial bounds exist. <L[2024-09-05-SimplexEfficiency] 10><L[2024-09-05-SimplexEfficiency] 11>%
\item \href{https://www.math.purdue.edu/~yipn/421/2024-09-05-SimplexEfficiency.pdf#page=12}{\textbf{Comparison of Simplex Method with Polynomial{-}Time Algorithms}}: Comparing the Simplex method with other linear programming algorithms that have proven polynomial-time complexity, such as interior-point methods. <L[2024-09-05-SimplexEfficiency] 12>%
\end{itemize}%
\begin{itemize}[[leftmargin=*]]%
\item \href{https://www.math.purdue.edu/~yipn/421/2024-09-12-Duality.pdf#page=1}{\textbf{Formulating Primal and Dual Linear Programs}}: This problem involves expressing a given optimization problem in both primal and dual forms, identifying the objective functions, constraints, and variables in each. <L[2024-09-12-Duality] 1><L[2024-09-12-Duality] 2><L[2024-09-12-Duality] 3><L[2024-09-12-Duality] 4>%
\item \href{https://www.math.purdue.edu/~yipn/421/2024-09-12-Duality.pdf#page=2}{\textbf{Applying Weak Duality Theorem}}: This problem involves verifying the inequality relationship between the objective function values of feasible primal and dual solutions. <L[2024-09-12-Duality] 2>%
\item \href{https://www.math.purdue.edu/~yipn/421/2024-09-12-Duality.pdf#page=3}{\textbf{Representing Linear Programs in Matrix Form}}: This problem focuses on translating primal and dual linear programs into their matrix representations, using matrices and vectors to represent the coefficients and variables. <L[2024-09-12-Duality] 3><L[2024-09-12-Duality] 4>%
\item \href{https://www.math.purdue.edu/~yipn/421/2024-09-12-Duality.pdf#page=5}{\textbf{Understanding Complementary Variables and their Relationship}}: This problem involves identifying and interpreting the relationships between primal and dual variables and their corresponding slack variables, particularly in the context of simplex iterations. <L[2024-09-12-Duality] 5>%
\item \href{https://www.math.purdue.edu/~yipn/421/2024-09-12-Duality.pdf#page=6}{\textbf{Demonstrating the Negative Transpose Property in Simplex Iterations}}: This problem involves showing how the negative transpose property is maintained throughout the simplex method iterations for both primal and dual problems. <L[2024-09-12-Duality] 6><L[2024-09-12-Duality] 7><L[2024-09-12-Duality] 8>%
\item \href{https://www.math.purdue.edu/~yipn/421/2024-09-12-Duality.pdf#page=9}{\textbf{Analyzing the Optimal Tableau and its Properties}}: This problem involves interpreting the characteristics of the optimal tableau for both primal and dual problems, specifically focusing on the non-negativity of coefficients. <L[2024-09-12-Duality] 9><L[2024-09-12-Duality] 10>%
\item \href{https://www.math.purdue.edu/~yipn/421/2024-09-12-Duality.pdf#page=11}{\textbf{Proving the Strong Duality Theorem}}: This problem involves demonstrating the equality of optimal objective function values for primal and dual problems, given the existence of an optimal solution for the primal problem. <L[2024-09-12-Duality] 11><L[2024-09-12-Duality] 12><L[2024-09-12-Duality] 13><L[2024-09-12-Duality] 14><L[2024-09-12-Duality] 15><L[2024-09-12-Duality] 16><L[2024-09-12-Duality] 17>%
\item \href{https://www.math.purdue.edu/~yipn/421/2024-09-12-Duality.pdf#page=18}{\textbf{Determining the Four Possible Outcomes of Primal{-}Dual Problem Pairs}}: This problem involves analyzing the feasibility and boundedness of primal and dual problems and determining which of the four possible outcomes applies to a given pair. <L[2024-09-12-Duality] 18><L[2024-09-12-Duality] 19><L[2024-09-12-Duality] 20>%
\item \href{https://www.math.purdue.edu/~yipn/421/2024-09-12-Duality.pdf#page=21}{\textbf{Applying the Complementary Slackness Theorem}}: This problem involves using the complementary slackness conditions to verify the optimality of feasible primal and dual solutions. <L[2024-09-12-Duality] 21><L[2024-09-12-Duality] 22><L[2024-09-12-Duality] 23>%
\end{itemize}%
\begin{itemize}[[leftmargin=*]]%
\item \href{https://www.math.purdue.edu/~yipn/421/2024-09-17-DualityGeneral.pdf#page=2}{\textbf{Formulating the Lagrangian for Equality Constraints}}: This problem involves constructing the Lagrangian function to incorporate equality constraints into the optimization problem using Lagrange multipliers. <L[2024-09-17-DualityGeneral] 2>%
\item \href{https://www.math.purdue.edu/~yipn/421/2024-09-17-DualityGeneral.pdf#page=7}{\textbf{Formulating the Lagrangian for Inequality and Equality Constraints}}: This problem extends the previous one by incorporating both inequality and equality constraints into the Lagrangian function, using separate Lagrange multipliers for each type of constraint. <L[2024-09-17-DualityGeneral] 7>%
\item \href{https://www.math.purdue.edu/~yipn/421/2024-09-17-DualityGeneral.pdf#page=5}{\textbf{Establishing Weak Duality}}: This problem focuses on proving that the optimal value of the primal problem is always less than or equal to the optimal value of the dual problem. <L[2024-09-17-DualityGeneral] 5><L[2024-09-17-DualityGeneral] 9>%
\item \href{https://www.math.purdue.edu/~yipn/421/2024-09-17-DualityGeneral.pdf#page=17}{\textbf{Establishing Strong Duality}}: This problem investigates the conditions under which the optimal values of the primal and dual problems are equal. <L[2024-09-17-DualityGeneral] 17><L[2024-09-17-DualityGeneral] 21><L[2024-09-17-DualityGeneral] 22>%
\item \href{https://www.math.purdue.edu/~yipn/421/2024-09-17-DualityGeneral.pdf#page=15}{\textbf{Analyzing the Maximum of the Lagrangian}}: This problem involves determining the maximum value of the Lagrangian function with respect to the primal variables, which is crucial for deriving the dual problem and establishing duality theorems. <L[2024-09-17-DualityGeneral] 15><L[2024-09-17-DualityGeneral] 16><L[2024-09-17-DualityGeneral] 20>%
\item \href{https://www.math.purdue.edu/~yipn/421/2024-09-17-DualityGeneral.pdf#page=11}{\textbf{Deriving the Dual Problem from the Lagrangian}}: This problem focuses on deriving the dual problem from the Lagrangian function by analyzing its maximum value with respect to the primal variables. <L[2024-09-17-DualityGeneral] 11><L[2024-09-17-DualityGeneral] 12><L[2024-09-17-DualityGeneral] 13><L[2024-09-17-DualityGeneral] 14>%
\item \href{https://www.math.purdue.edu/~yipn/421/2024-09-17-DualityGeneral.pdf#page=11}{\textbf{Applying Duality to a General Linear Program}}: This problem involves formulating and analyzing the primal and dual problems for a general linear program in standard form. <L[2024-09-17-DualityGeneral] 11><L[2024-09-17-DualityGeneral] 12><L[2024-09-17-DualityGeneral] 13><L[2024-09-17-DualityGeneral] 14><L[2024-09-17-DualityGeneral] 15><L[2024-09-17-DualityGeneral] 16><L[2024-09-17-DualityGeneral] 17><L[2024-09-17-DualityGeneral] 18>%
\item \href{https://www.math.purdue.edu/~yipn/421/2024-09-17-DualityGeneral.pdf#page=19}{\textbf{Applying Duality to a "Traditional" Linear Program}}: This problem focuses on applying duality concepts to the specific form of linear programs often encountered in applications. <L[2024-09-17-DualityGeneral] 19><L[2024-09-17-DualityGeneral] 20><L[2024-09-17-DualityGeneral] 21><L[2024-09-17-DualityGeneral] 22>%
\item \href{https://www.math.purdue.edu/~yipn/421/2024-09-17-DualityGeneral.pdf#page=22}{\textbf{Understanding Complementary Slackness}}: This problem involves understanding and applying the complementary slackness conditions, which relate the optimal solutions of the primal and dual problems. <L[2024-09-17-DualityGeneral] 22>%
\end{itemize}%
\begin{itemize}[[leftmargin=*]]%
\item \href{https://www.math.purdue.edu/~yipn/421/2024-09-19-DualityExamples.pdf#page=5}{\textbf{Formulating the Diet Problem as a Linear Program}}: This problem involves translating the description of the diet problem (minimizing cost while meeting nutritional requirements) into a mathematical model with an objective function and constraints. <L[2024-09-19-DualityExamples] 5>%
\item \href{https://www.math.purdue.edu/~yipn/421/2024-09-19-DualityExamples.pdf#page=8}{\textbf{Constructing the Dual of the Diet Problem}}: This problem focuses on formulating the dual linear program corresponding to the primal diet problem, which involves maximizing profit from selling nutrient pills. <L[2024-09-19-DualityExamples] 8>%
\item \href{https://www.math.purdue.edu/~yipn/421/2024-09-19-DualityExamples.pdf#page=3}{\textbf{Representing the Diet Problem using a Table}}: This problem involves representing the nutritional content of foods and their costs in a tabular format to facilitate the formulation of the linear program. <L[2024-09-19-DualityExamples] 3><L[2024-09-19-DualityExamples] 4><L[2024-09-19-DualityExamples] 6>%
\item \href{https://www.math.purdue.edu/~yipn/421/2024-09-19-DualityExamples.pdf#page=1}{\textbf{Understanding the Relationship Between Primal and Dual Problems in the Context of the Diet Problem}}: This problem involves interpreting the primal (minimizing cost of food) and dual (maximizing profit from pills) problems and understanding their connection through the duality theorem. <L[2024-09-19-DualityExamples] 1><L[2024-09-19-DualityExamples] 7><L[2024-09-19-DualityExamples] 8>%
\end{itemize}%
\begin{itemize}[[leftmargin=*]]%
\item \href{https://www.math.purdue.edu/~yipn/421/2024-09-24-SimplexMatrix.pdf#page=1}{\textbf{Formulating a Linear Program in Matrix Notation}}: Representing a linear program using matrices and vectors, including the objective function, constraints, and non-negativity conditions. <L[2024-09-24-SimplexMatrix] 1><L[2024-09-24-SimplexMatrix] 2><L[2024-09-24-SimplexMatrix] 3>%
\item \href{https://www.math.purdue.edu/~yipn/421/2024-09-24-SimplexMatrix.pdf#page=2}{\textbf{Transforming the Linear Program for Simplex Iterations}}: Modifying the matrix representation of a linear program to a form suitable for iterative solution using the simplex method. <L[2024-09-24-SimplexMatrix] 2><L[2024-09-24-SimplexMatrix] 3>%
\item \href{https://www.math.purdue.edu/~yipn/421/2024-09-24-SimplexMatrix.pdf#page=4}{\textbf{Partitioning Variables into Basic and Non{-}Basic Sets}}: Dividing the variables of a linear program into basic and non-basic sets to facilitate the simplex method's iterative process. <L[2024-09-24-SimplexMatrix] 4><L[2024-09-24-SimplexMatrix] 5><L[2024-09-24-SimplexMatrix] 6><L[2024-09-24-SimplexMatrix] 7>%
\item \href{https://www.math.purdue.edu/~yipn/421/2024-09-24-SimplexMatrix.pdf#page=8}{\textbf{Expressing the Objective Function in Terms of Basic and Non{-}Basic Variables}}: Rewriting the objective function using the basic and non-basic variable partitions to enable efficient updates during simplex iterations. <L[2024-09-24-SimplexMatrix] 8><L[2024-09-24-SimplexMatrix] 9>%
\item \href{https://www.math.purdue.edu/~yipn/421/2024-09-24-SimplexMatrix.pdf#page=10}{\textbf{Updating the Simplex Tableau Using Matrix Operations}}: Performing matrix calculations to update the simplex tableau during each iteration of the simplex method. <L[2024-09-24-SimplexMatrix] 10><L[2024-09-24-SimplexMatrix] 11><L[2024-09-24-SimplexMatrix] 12><L[2024-09-24-SimplexMatrix] 13>%
\item \href{https://www.math.purdue.edu/~yipn/421/2024-09-24-SimplexMatrix.pdf#page=14}{\textbf{Identifying Complementary Relationships Between Primal and Dual Variables}}: Recognizing the complementary slackness conditions between primal and dual variables in the simplex method. <L[2024-09-24-SimplexMatrix] 14>%
\item \href{https://www.math.purdue.edu/~yipn/421/2024-09-24-SimplexMatrix.pdf#page=16}{\textbf{Solving a Linear Program Graphically}}: Determining the optimal solution of a linear program by visualizing the feasible region and objective function. <L[2024-09-24-SimplexMatrix] 16><L[2024-09-24-SimplexMatrix] 21>%
\item \href{https://www.math.purdue.edu/~yipn/421/2024-09-24-SimplexMatrix.pdf#page=15}{\textbf{Iteratively Solving a Linear Program Using the Simplex Method}}: Applying the simplex method to find the optimal solution of a linear program through a series of iterations. <L[2024-09-24-SimplexMatrix] 15><L[2024-09-24-SimplexMatrix] 17><L[2024-09-24-SimplexMatrix] 18><L[2024-09-24-SimplexMatrix] 19><L[2024-09-24-SimplexMatrix] 20>%
\end{itemize}%
\begin{itemize}[[leftmargin=*]]%
\item \href{https://www.math.purdue.edu/~yipn/421/2024-09-26-Sensitivity.pdf#page=4}{\textbf{Determining the range of objective function coefficient changes that maintain optimality}}: This problem investigates how much the objective function coefficients can change before the optimal solution changes. The analysis focuses on maintaining non-negativity of the reduced costs of non-basic variables. <L[2024-09-26-Sensitivity] 4><L[2024-09-26-Sensitivity] 5><L[2024-09-26-Sensitivity] 6><L[2024-09-26-Sensitivity] 7>%
\item \href{https://www.math.purdue.edu/~yipn/421/2024-09-26-Sensitivity.pdf#page=8}{\textbf{Determining the range of constraint value changes that maintain optimality}}: This problem investigates how much the right-hand side values of the constraints can change before the optimal solution changes. The analysis focuses on maintaining non-negativity of the basic variables. <L[2024-09-26-Sensitivity] 8><L[2024-09-26-Sensitivity] 9><L[2024-09-26-Sensitivity] 10><L[2024-09-26-Sensitivity] 11>%
\end{itemize}%
\begin{itemize}[[leftmargin=*]]%
\item \href{https://www.math.purdue.edu/~yipn/421/2024-10-01-DualGeneralLP.pdf#page=1}{\textbf{Deriving the Dual of a General Linear Program}}: This problem focuses on deriving the dual linear program from a given primal linear program, particularly when the primal problem does not have non-negativity constraints on all variables. The derivation involves techniques like introducing slack variables and using the Lagrangian function. <L[2024-10-01-DualGeneralLP] 1><L[2024-10-01-DualGeneralLP] 2><L[2024-10-01-DualGeneralLP] 3><L[2024-10-01-DualGeneralLP] 4><L[2024-10-01-DualGeneralLP] 5><L[2024-10-01-DualGeneralLP] 6><L[2024-10-01-DualGeneralLP] 7>%
\item \href{https://www.math.purdue.edu/~yipn/421/2024-10-01-DualGeneralLP.pdf#page=5}{\textbf{Proving Weak Duality}}: This problem involves demonstrating that the objective function value of any feasible solution to the primal problem is less than or equal to the objective function value of any feasible solution to the dual problem. <L[2024-10-01-DualGeneralLP] 5>%
\end{itemize}%
\begin{itemize}[[leftmargin=*]]%
\item \href{https://www.math.purdue.edu/~yipn/421/2024-10-01-SensitivityDuality.pdf#page=13}{\textbf{Verifying Optimality}}: A method for verifying the optimality of a solution is presented, involving checking the objective function value and the feasibility of the dual solution. <L[2024-10-01-SensitivityDuality] 13>%
\item \textbf{Sensitivity Analysis}: Determining how changes in constraint values (specifically the right-hand side) affect the optimal solution and objective function value. <SLIDE 1, SLIDE 7, SLIDE 8, SLIDE 9, SLIDE 10>%
\item \href{https://www.math.purdue.edu/~yipn/421/2024-10-01-SensitivityDuality.pdf#page=3}{\textbf{Applying Complementary Slackness}}: Using complementary slackness conditions to verify the optimal solution by relating primal and dual variables. <L[2024-10-01-SensitivityDuality] 3>%
\item \textbf{Solving Linear Programs using the Basis and Non{-}Basis Method}: Utilizing the basis and non-basis method (likely the simplex method) to solve linear programming problems, involving matrix operations and iterative calculations. <SLIDE 4, SLIDE 5, SLIDE 6, SLIDE 7, SLIDE 11, SLIDE 12, SLIDE 14, SLIDE 15>%
\item \textbf{Understanding Duality}: Exploring the relationship between the primal and dual linear programming problems, including strong duality and complementary slackness. <SLIDE 2, SLIDE 3, SLIDE 11, SLIDE 12, SLIDE 14, SLIDE 15>%
\item \href{https://www.math.purdue.edu/~yipn/421/2024-10-01-SensitivityDuality.pdf#page=10}{\textbf{Determining Shadow Prices}}: Calculating the shadow prices (dual variables) to determine the rate of change in the optimal objective function value for a unit change in the right-hand side of a constraint. <L[2024-10-01-SensitivityDuality] 10>%
\end{itemize}%
\begin{itemize}[[leftmargin=*]]%
\item \href{https://www.math.purdue.edu/~yipn/421/2024-10-03-DualSimplexMatrix.pdf#page=1}{\textbf{Solving Primal and Dual Problems Simultaneously}}: This problem focuses on finding the optimal solutions for both the primal and dual linear programming problems, highlighting their relationship and the equality of their optimal objective function values. <L[2024-10-03-DualSimplexMatrix] 1>%
\item \href{https://www.math.purdue.edu/~yipn/421/2024-10-03-DualSimplexMatrix.pdf#page=2}{\textbf{Formulating and Solving Primal and Dual Linear Programs}}: This problem involves constructing the primal and dual linear programming formulations from a given problem statement and then solving both using appropriate methods like the dual simplex method. <L[2024-10-03-DualSimplexMatrix] 2><L[2024-10-03-DualSimplexMatrix] 3><L[2024-10-03-DualSimplexMatrix] 4><L[2024-10-03-DualSimplexMatrix] 5><L[2024-10-03-DualSimplexMatrix] 6>%
\item \href{https://www.math.purdue.edu/~yipn/421/2024-10-03-DualSimplexMatrix.pdf#page=2}{\textbf{Applying the Dual Simplex Method}}: This problem involves using the dual simplex method to iteratively solve a linear programming problem, starting from a dual-feasible but primal-infeasible solution. <L[2024-10-03-DualSimplexMatrix] 2><L[2024-10-03-DualSimplexMatrix] 3><L[2024-10-03-DualSimplexMatrix] 4>%
\item \href{https://www.math.purdue.edu/~yipn/421/2024-10-03-DualSimplexMatrix.pdf#page=5}{\textbf{Implementing a Dual Based Phase I Algorithm}}: This problem focuses on using a Phase I algorithm that leverages the dual problem to find a feasible starting point for the simplex method when the origin is infeasible for both the primal and dual problems. <L[2024-10-03-DualSimplexMatrix] 5><L[2024-10-03-DualSimplexMatrix] 6>%
\item \href{https://www.math.purdue.edu/~yipn/421/2024-10-03-DualSimplexMatrix.pdf#page=7}{\textbf{Expressing the Simplex Method in Matrix Form}}: This problem involves representing the simplex method using matrix notation, including the objective function, constraints, and the update rules for the basis matrices. <L[2024-10-03-DualSimplexMatrix] 7><L[2024-10-03-DualSimplexMatrix] 8><L[2024-10-03-DualSimplexMatrix] 9><L[2024-10-03-DualSimplexMatrix] 10><L[2024-10-03-DualSimplexMatrix] 11><L[2024-10-03-DualSimplexMatrix] 12><L[2024-10-03-DualSimplexMatrix] 13><L[2024-10-03-DualSimplexMatrix] 14><L[2024-10-03-DualSimplexMatrix] 15>%
\item \href{https://www.math.purdue.edu/~yipn/421/2024-10-03-DualSimplexMatrix.pdf#page=14}{\textbf{Efficiently Updating Basis Matrices in the Simplex Method}}: This problem focuses on developing efficient methods for updating the basis matrices and their inverses during the simplex iterations, using elementary matrices or other techniques to avoid redundant computations. <L[2024-10-03-DualSimplexMatrix] 14><L[2024-10-03-DualSimplexMatrix] 15>%
\end{itemize}%
\begin{itemize}[[leftmargin=*]]%
\item \href{https://www.math.purdue.edu/~yipn/421/2024-10-22-Convexity.pdf#page=1}{\textbf{Defining Convex Sets and Convex Combinations}}: This problem involves formally defining convex sets and convex combinations, and illustrating these concepts geometrically in $\mathbb{R}^2$. <L[2024-10-22-Convexity] 1>%
\item \href{https://www.math.purdue.edu/~yipn/421/2024-10-22-Convexity.pdf#page=2}{\textbf{Proving the Equivalence of Convexity and Convex Combinations}}: This problem focuses on proving that a set is convex if and only if it contains all convex combinations of its points. <L[2024-10-22-Convexity] 2><L[2024-10-22-Convexity] 3><L[2024-10-22-Convexity] 4>%
\item \href{https://www.math.purdue.edu/~yipn/421/2024-10-22-Convexity.pdf#page=5}{\textbf{Defining and Proving Properties of Convex Hulls}}: This problem involves defining the convex hull of a set and proving that it's equivalent to the set of all convex combinations of finitely many points from the original set. <L[2024-10-22-Convexity] 5><L[2024-10-22-Convexity] 6><L[2024-10-22-Convexity] 7><L[2024-10-22-Convexity] 8>%
\item \href{https://www.math.purdue.edu/~yipn/421/2024-10-22-Convexity.pdf#page=9}{\textbf{Proving the Caratheodory Theorem}}: This problem involves proving that any point in the convex hull of a set in $\mathbb{R}^m$ can be expressed as a convex combination of at most $m+1$ points from the original set, using linear programming. <L[2024-10-22-Convexity] 9>%
\item \href{https://www.math.purdue.edu/~yipn/421/2024-10-22-Convexity.pdf#page=10}{\textbf{Proving the Separation Theorem}}: This problem focuses on proving that two disjoint nonempty polyhedra in $\mathbb{R}^n$ can be separated by two disjoint half-spaces. The proof utilizes Farkas' Lemma. <L[2024-10-22-Convexity] 10><L[2024-10-22-Convexity] 11><L[2024-10-22-Convexity] 12><L[2024-10-22-Convexity] 13><L[2024-10-22-Convexity] 14><L[2024-10-22-Convexity] 15>%
\item \href{https://www.math.purdue.edu/~yipn/421/2024-10-22-Convexity.pdf#page=15}{\textbf{Proving Farkas' Lemma}}: This problem involves proving Farkas' Lemma, a fundamental result in linear programming, which provides a condition for the infeasibility of a system of linear inequalities. <L[2024-10-22-Convexity] 15>%
\end{itemize}%
\begin{itemize}[[leftmargin=*]]%
\item \href{https://www.math.purdue.edu/~yipn/421/2024-10-22-Regression.pdf#page=4}{\textbf{Linear Regression with L² Error Minimization}}: Finding parameters that minimize the sum of squared errors between observed and predicted values in a linear model. <L[2024-10-22-Regression] 4>%
\item \href{https://www.math.purdue.edu/~yipn/421/2024-10-22-Regression.pdf#page=5}{\textbf{Linear Regression with L¹ Error Minimization}}: Finding parameters that minimize the sum of absolute errors between observed and predicted values in a linear model. <L[2024-10-22-Regression] 5>%
\item \href{https://www.math.purdue.edu/~yipn/421/2024-10-22-Regression.pdf#page=6}{\textbf{Linear Regression with L∞ Error Minimization}}: Finding parameters that minimize the maximum absolute error between observed and predicted values in a linear model. <L[2024-10-22-Regression] 6>%
\item \href{https://www.math.purdue.edu/~yipn/421/2024-10-22-Regression.pdf#page=7}{\textbf{Binary Classification with Linear Separability}}: Determining if a hyperplane exists that perfectly separates two classes of data points. <L[2024-10-22-Regression] 7><L[2024-10-22-Regression] 8>%
\item \href{https://www.math.purdue.edu/~yipn/421/2024-10-22-Regression.pdf#page=9}{\textbf{Maximum Margin Classifier}}: Finding the hyperplane that maximizes the minimum distance to the closest data points of each class. <L[2024-10-22-Regression] 9>%
\item \href{https://www.math.purdue.edu/~yipn/421/2024-10-22-Regression.pdf#page=10}{\textbf{Support Vector Machine (SVM)}}: Finding the hyperplane that maximizes the margin between two classes of data points. <L[2024-10-22-Regression] 10>%
\item \href{https://www.math.purdue.edu/~yipn/421/2024-10-22-Regression.pdf#page=11}{\textbf{SVM Optimization as a Quadratic Programming Problem}}: Formulating the SVM problem as a quadratic programming problem to find the optimal hyperplane. <L[2024-10-22-Regression] 11><L[2024-10-22-Regression] 13><L[2024-10-22-Regression] 14><L[2024-10-22-Regression] 19><L[2024-10-22-Regression] 20>%
\item \href{https://www.math.purdue.edu/~yipn/421/2024-10-22-Regression.pdf#page=15}{\textbf{Finding the Maximum Inscribed Circle in a Convex Polyhedron}}: Determining the largest circle that can fit entirely within a given convex polyhedron. <L[2024-10-22-Regression] 15><L[2024-10-22-Regression] 16><L[2024-10-22-Regression] 21>%
\item \href{https://www.math.purdue.edu/~yipn/421/2024-10-22-Regression.pdf#page=22}{\textbf{Smallest Enclosing Ball Problem}}: Finding the smallest radius ball that contains all given points in a d-dimensional space. <L[2024-10-22-Regression] 22>%
\end{itemize}%
\begin{itemize}[[leftmargin=*]]%
\item \href{https://www.math.purdue.edu/~yipn/421/2024-10-28-ApplsLinProg.pdf#page=3}{\textbf{Production Planning}}: Developing a yearly production plan for an ice cream manufacturer to meet fluctuating monthly demands while minimizing production change and storage costs. <L[2024-10-28-ApplsLinProg] 3><L[2024-10-28-ApplsLinProg] 4><L[2024-10-28-ApplsLinProg] 5>%
\item \href{https://www.math.purdue.edu/~yipn/421/2024-10-28-ApplsLinProg.pdf#page=7}{\textbf{Workforce Planning}}: Determining the optimal number of regular and temporary employees to minimize labor costs while meeting fluctuating monthly labor demands. <L[2024-10-28-ApplsLinProg] 7><L[2024-10-28-ApplsLinProg] 8>%
\item \href{https://www.math.purdue.edu/~yipn/421/2024-10-28-ApplsLinProg.pdf#page=9}{\textbf{Portfolio Selection (Markowitz Model)}}: Optimizing an investment portfolio to maximize expected return while minimizing risk, considering both average return and risk (variance). <L[2024-10-28-ApplsLinProg] 9><L[2024-10-28-ApplsLinProg] 10><L[2024-10-28-ApplsLinProg] 11><L[2024-10-28-ApplsLinProg] 12><L[2024-10-28-ApplsLinProg] 13><L[2024-10-28-ApplsLinProg] 14><L[2024-10-28-ApplsLinProg] 16>%
\item \href{https://www.math.purdue.edu/~yipn/421/2024-10-28-ApplsLinProg.pdf#page=17}{\textbf{Smallest Enclosing Ball Problem}}: Finding the smallest ball (hypersphere) that contains a given set of points in a d-dimensional space. <L[2024-10-28-ApplsLinProg] 17><L[2024-10-28-ApplsLinProg] 18>%
\item \href{https://www.math.purdue.edu/~yipn/421/2024-10-28-ApplsLinProg.pdf#page=19}{\textbf{Maximum Weight Matching}}: Finding a matching in a weighted graph that maximizes the total weight of the selected edges. <L[2024-10-28-ApplsLinProg] 19><L[2024-10-28-ApplsLinProg] 20><L[2024-10-28-ApplsLinProg] 21>%
\item \href{https://www.math.purdue.edu/~yipn/421/2024-10-28-ApplsLinProg.pdf#page=22}{\textbf{Machine Scheduling}}: Determining the optimal order to process printing jobs on different machines to minimize the total processing time. <L[2024-10-28-ApplsLinProg] 22><L[2024-10-28-ApplsLinProg] 23>%
\item \href{https://www.math.purdue.edu/~yipn/421/2024-10-28-ApplsLinProg.pdf#page=24}{\textbf{Knapsack Problem}}: Selecting a subset of items with weights and values to maximize the total value while not exceeding a weight capacity. <L[2024-10-28-ApplsLinProg] 24>%
\item \href{https://www.math.purdue.edu/~yipn/421/2024-10-28-ApplsLinProg.pdf#page=25}{\textbf{Cutting Stock Problem}}: Determining how to cut large rolls of paper into smaller rolls of various widths to minimize waste while meeting customer demands. <L[2024-10-28-ApplsLinProg] 25><L[2024-10-28-ApplsLinProg] 26><L[2024-10-28-ApplsLinProg] 27><L[2024-10-28-ApplsLinProg] 28><L[2024-10-28-ApplsLinProg] 29>%
\item \href{https://www.math.purdue.edu/~yipn/421/2024-10-28-ApplsLinProg.pdf#page=30}{\textbf{Sparse Solutions of Linear System}}: Finding a solution to an underdetermined system of linear equations that has a limited number of non-zero elements. <L[2024-10-28-ApplsLinProg] 30><L[2024-10-28-ApplsLinProg] 31><L[2024-10-28-ApplsLinProg] 32><L[2024-10-28-ApplsLinProg] 33><L[2024-10-28-ApplsLinProg] 34><L[2024-10-28-ApplsLinProg] 35><L[2024-10-28-ApplsLinProg] 36><L[2024-10-28-ApplsLinProg] 37>%
\end{itemize}%
\begin{itemize}[[leftmargin=*]]%
\item \href{https://www.math.purdue.edu/~yipn/421/2024-10-28-Farkas.pdf#page=2}{\textbf{Determining the existence of a non{-}negative solution to `Ax = b`}}: This problem investigates whether there exists a vector `x` with non-negative components that satisfies the equation `Ax = b`. Several equivalent conditions are presented using a vector `y`. <L[2024-10-28-Farkas] 2><L[2024-10-28-Farkas] 3><L[2024-10-28-Farkas] 4><L[2024-10-28-Farkas] 5><L[2024-10-28-Farkas] 6><L[2024-10-28-Farkas] 7><L[2024-10-28-Farkas] 8><L[2024-10-28-Farkas] 9><L[2024-10-28-Farkas] 10><L[2024-10-28-Farkas] 11>%
\item \href{https://www.math.purdue.edu/~yipn/421/2024-10-28-Farkas.pdf#page=4}{\textbf{Determining the existence of a solution to `Ax ≤ b`}}: This problem explores whether a solution `x` exists that satisfies the system of inequalities `Ax ≤ b`. Different versions explore the existence of non-negative solutions. <L[2024-10-28-Farkas] 4><L[2024-10-28-Farkas] 5><L[2024-10-28-Farkas] 6><L[2024-10-28-Farkas] 7><L[2024-10-28-Farkas] 8><L[2024-10-28-Farkas] 9><L[2024-10-28-Farkas] 10><L[2024-10-28-Farkas] 11>%
\item \href{https://www.math.purdue.edu/~yipn/421/2024-10-28-Farkas.pdf#page=1}{\textbf{Proving Farkas' Lemma}}: This problem focuses on proving the different variants of Farkas' Lemma, which provides conditions for the existence or non-existence of solutions to systems of linear inequalities and equations of the form `Ax = b` and `Ax ≤ b`, with or without non-negativity constraints on `x`. <L[2024-10-28-Farkas] 1><L[2024-10-28-Farkas] 2><L[2024-10-28-Farkas] 3><L[2024-10-28-Farkas] 4><L[2024-10-28-Farkas] 5><L[2024-10-28-Farkas] 6><L[2024-10-28-Farkas] 7><L[2024-10-28-Farkas] 8><L[2024-10-28-Farkas] 9><L[2024-10-28-Farkas] 10><L[2024-10-28-Farkas] 11>%
\end{itemize}%
\begin{itemize}[[leftmargin=*]]%
\item \href{https://www.math.purdue.edu/~yipn/421/2024-10-31-IntLinProg.pdf#page=2}{\textbf{Solving Integer Programming Problems by Relaxing and Rounding}}: This problem explores the limitations of solving the relaxed version of an integer program and then rounding the solution to obtain an integer solution. It highlights that the rounded solution may be infeasible or far from the optimal integer solution. <L[2024-10-31-IntLinProg] 2>%
\item \href{https://www.math.purdue.edu/~yipn/421/2024-10-31-IntLinProg.pdf#page=3}{\textbf{Solving a Specific Integer Programming Problem}}: This problem involves finding the optimal integer solution to a given integer programming problem with a linear objective function and linear constraints. A graphical method and the Branch and Bound method are used to solve this problem. <L[2024-10-31-IntLinProg] 3><L[2024-10-31-IntLinProg] 4><L[2024-10-31-IntLinProg] 5><L[2024-10-31-IntLinProg] 6><L[2024-10-31-IntLinProg] 7><L[2024-10-31-IntLinProg] 8><L[2024-10-31-IntLinProg] 9><L[2024-10-31-IntLinProg] 10>%
\item \href{https://www.math.purdue.edu/~yipn/421/2024-10-31-IntLinProg.pdf#page=5}{\textbf{Applying the Branch and Bound Method}}: This problem focuses on using the Branch and Bound algorithm to solve an integer programming problem. It involves creating a branching tree, solving relaxed linear programs at each node, and pruning branches based on the best integer solution found so far. Depth-first search is discussed as an efficient strategy for exploring the tree. <L[2024-10-31-IntLinProg] 5><L[2024-10-31-IntLinProg] 6><L[2024-10-31-IntLinProg] 7><L[2024-10-31-IntLinProg] 8><L[2024-10-31-IntLinProg] 9><L[2024-10-31-IntLinProg] 10><L[2024-10-31-IntLinProg] 11><L[2024-10-31-IntLinProg] 12><L[2024-10-31-IntLinProg] 13><L[2024-10-31-IntLinProg] 14>%
\item \href{https://www.math.purdue.edu/~yipn/421/2024-10-31-IntLinProg.pdf#page=15}{\textbf{Applying Gomory Cuts to Solve Integer Programming Problems}}: This problem demonstrates the use of Gomory cuts to find an integer solution to an integer programming problem. It involves solving the relaxed problem, identifying fractional parts in the optimal dictionary, generating Gomory cuts, and iteratively adding these cuts to the problem until an integer solution is found. <L[2024-10-31-IntLinProg] 15><L[2024-10-31-IntLinProg] 16><L[2024-10-31-IntLinProg] 17><L[2024-10-31-IntLinProg] 18><L[2024-10-31-IntLinProg] 19><L[2024-10-31-IntLinProg] 20><L[2024-10-31-IntLinProg] 21><L[2024-10-31-IntLinProg] 22><L[2024-10-31-IntLinProg] 23>%
\end{itemize}%
\begin{itemize}[[leftmargin=*]]%
\item \href{https://www.math.purdue.edu/~yipn/421/2024-11-05-NetFlows.pdf#page=2}{\textbf{Formulating a Network Flow Problem as a Linear Program}}: This problem involves representing a network flow problem using the standard linear programming framework, defining the objective function (minimizing total cost) and constraints (flow conservation at each node and non-negativity of flows). <L[2024-11-05-NetFlows] 2><L[2024-11-05-NetFlows] 3><L[2024-11-05-NetFlows] 4><L[2024-11-05-NetFlows] 7><L[2024-11-05-NetFlows] 8><L[2024-11-05-NetFlows] 9>%
\item \href{https://www.math.purdue.edu/~yipn/421/2024-11-05-NetFlows.pdf#page=5}{\textbf{Analyzing the Incidence Matrix of a Network}}: This problem focuses on understanding the properties of the incidence matrix, including its rank and the relationship between its submatrices and spanning trees of the network. <L[2024-11-05-NetFlows] 5><L[2024-11-05-NetFlows] 6><L[2024-11-05-NetFlows] 15><L[2024-11-05-NetFlows] 16><L[2024-11-05-NetFlows] 17>%
\item \href{https://www.math.purdue.edu/~yipn/421/2024-11-05-NetFlows.pdf#page=10}{\textbf{Finding a Spanning Tree Solution}}: This problem involves finding a feasible solution to the network flow problem where the flow along arcs in a spanning tree is zero. This solution might not be optimal. <L[2024-11-05-NetFlows] 10><L[2024-11-05-NetFlows] 11><L[2024-11-05-NetFlows] 12><L[2024-11-05-NetFlows] 13><L[2024-11-05-NetFlows] 14>%
\item \href{https://www.math.purdue.edu/~yipn/421/2024-11-05-NetFlows.pdf#page=15}{\textbf{Relating Spanning Trees to Bases in the Simplex Method}}: This problem explores the connection between spanning trees in a network and bases in the simplex method for solving linear programs. It involves understanding how spanning trees can be used to construct bases for the incidence matrix. <L[2024-11-05-NetFlows] 15><L[2024-11-05-NetFlows] 16><L[2024-11-05-NetFlows] 17>%
\item \href{https://www.math.purdue.edu/~yipn/421/2024-11-05-NetFlows.pdf#page=22}{\textbf{Applying the Primal Simplex Method to Network Flows}}: This problem involves using the primal simplex method to iteratively improve a spanning tree solution until an optimal solution is found. This involves identifying entering and leaving variables, updating the spanning tree, and updating primal and dual variables. <L[2024-11-05-NetFlows] 22>%
\item \href{https://www.math.purdue.edu/~yipn/421/2024-11-05-NetFlows.pdf#page=18}{\textbf{Determining Optimality using Complementary Slackness}}: This problem involves verifying the optimality of a network flow solution by checking the complementary slackness conditions, which relate primal flows and dual slack variables. <L[2024-11-05-NetFlows] 18><L[2024-11-05-NetFlows] 19><L[2024-11-05-NetFlows] 20><L[2024-11-05-NetFlows] 21>%
\end{itemize}%
\begin{itemize}[[leftmargin=*]]%
\item \href{https://www.math.purdue.edu/~yipn/421/2024-11-07-NetSimplex.pdf#page=2}{\textbf{Finding a Primal Feasible Flow}}: Determining a flow that satisfies the constraints $\tilde{B}\tilde{X} = \tilde{b}$ and $x_{ij} = 0$, $(i,j) \in A \setminus T$, potentially involving negative flow values. <L[2024-11-07-NetSimplex] 2>%
\item \href{https://www.math.purdue.edu/~yipn/421/2024-11-07-NetSimplex.pdf#page=4}{\textbf{Modifying the Spanning Tree}}: Understanding how adding or removing an arc from a spanning tree creates a cycle or disconnects the tree, respectively. <L[2024-11-07-NetSimplex] 4>%
\item \href{https://www.math.purdue.edu/~yipn/421/2024-11-07-NetSimplex.pdf#page=6}{\textbf{Primal Pivot Step}}: Selecting an entering arc with $z_{ij} < 0$, finding a cycle, and updating primal and dual flows to improve the solution. <L[2024-11-07-NetSimplex] 6><L[2024-11-07-NetSimplex] 7><L[2024-11-07-NetSimplex] 8><L[2024-11-07-NetSimplex] 9><L[2024-11-07-NetSimplex] 10><L[2024-11-07-NetSimplex] 11><L[2024-11-07-NetSimplex] 12><L[2024-11-07-NetSimplex] 13><L[2024-11-07-NetSimplex] 14><L[2024-11-07-NetSimplex] 15><L[2024-11-07-NetSimplex] 16><L[2024-11-07-NetSimplex] 17><L[2024-11-07-NetSimplex] 18><L[2024-11-07-NetSimplex] 19><L[2024-11-07-NetSimplex] 20><L[2024-11-07-NetSimplex] 21><L[2024-11-07-NetSimplex] 22><L[2024-11-07-NetSimplex] 23><L[2024-11-07-NetSimplex] 24><L[2024-11-07-NetSimplex] 25><L[2024-11-07-NetSimplex] 26><L[2024-11-07-NetSimplex] 27><L[2024-11-07-NetSimplex] 28><L[2024-11-07-NetSimplex] 30><L[2024-11-07-NetSimplex] 31>%
\item \href{https://www.math.purdue.edu/~yipn/421/2024-11-07-NetSimplex.pdf#page=32}{\textbf{Dual Pivot Step}}: Selecting a leaving arc with a negative primal flow, finding two subtrees, and selecting an entering arc to reconnect them while maintaining dual feasibility. <L[2024-11-07-NetSimplex] 32><L[2024-11-07-NetSimplex] 34><L[2024-11-07-NetSimplex] 35><L[2024-11-07-NetSimplex] 36><L[2024-11-07-NetSimplex] 37><L[2024-11-07-NetSimplex] 38><L[2024-11-07-NetSimplex] 39><L[2024-11-07-NetSimplex] 40><L[2024-11-07-NetSimplex] 41><L[2024-11-07-NetSimplex] 42>%
\item \href{https://www.math.purdue.edu/~yipn/421/2024-11-07-NetSimplex.pdf#page=45}{\textbf{Developing Two{-}Phased Network Simplex Methods}}: Creating methods that first find a primal or dual feasible solution using artificial costs or supplies, then optimize using the primal or dual network simplex method. <L[2024-11-07-NetSimplex] 45>%
\item \href{https://www.math.purdue.edu/~yipn/421/2024-11-07-NetSimplex.pdf#page=45}{\textbf{Developing Parametric Self{-}Dual Network Simplex Methods}}: Creating methods that interweave primal and dual pivots to reduce a perturbation parameter from infinity to zero. <L[2024-11-07-NetSimplex] 45>%
\end{itemize}%
\begin{itemize}[[leftmargin=*]]%
\item \href{https://www.math.purdue.edu/~yipn/421/2024-11-14-NetAppls.pdf#page=1}{\textbf{Modeling Transportation Problems with Inequality Constraints}}: This problem involves formulating a linear program to model a transportation problem where supply and demand may not be perfectly balanced, requiring the use of inequality constraints and potentially a dummy node to represent excess supply or unmet demand. <L[2024-11-14-NetAppls] 1><L[2024-11-14-NetAppls] 2><L[2024-11-14-NetAppls] 3><L[2024-11-14-NetAppls] 4><L[2024-11-14-NetAppls] 5><L[2024-11-14-NetAppls] 6>%
\item \href{https://www.math.purdue.edu/~yipn/421/2024-11-14-NetAppls.pdf#page=3}{\textbf{Formulating Transportation Problems with Dummy Nodes}}: This problem focuses on using dummy nodes in the network representation of a transportation problem to handle situations where supply exceeds demand or vice versa, transforming inequality constraints into equality constraints. <L[2024-11-14-NetAppls] 3><L[2024-11-14-NetAppls] 4>%
\item \href{https://www.math.purdue.edu/~yipn/421/2024-11-14-NetAppls.pdf#page=7}{\textbf{Generalizing Inequality Constraints in Network Flows}}: This problem involves expressing general inequality constraints in network flow problems using matrices and vectors, converting inequalities to equalities by introducing slack variables. <L[2024-11-14-NetAppls] 7>%
\item \href{https://www.math.purdue.edu/~yipn/421/2024-11-14-NetAppls.pdf#page=9}{\textbf{Solving the Upper Bounded Transshipment Problem}}: This problem involves finding optimal flows in a network where each arc has an upper bound on the flow it can carry, requiring the use of additional constraints and potentially slack variables. <L[2024-11-14-NetAppls] 9><L[2024-11-14-NetAppls] 10><L[2024-11-14-NetAppls] 11><L[2024-11-14-NetAppls] 12>%
\item \href{https://www.math.purdue.edu/~yipn/421/2024-11-14-NetAppls.pdf#page=13}{\textbf{Solving the Shortest Path Problem using Network Flow}}: This problem involves formulating and solving the shortest path problem as a linear program using a network flow approach, with constraints ensuring flow conservation and a single unit of flow from source to destination. <L[2024-11-14-NetAppls] 13><L[2024-11-14-NetAppls] 14><L[2024-11-14-NetAppls] 15>%
\item \href{https://www.math.purdue.edu/~yipn/421/2024-11-14-NetAppls.pdf#page=16}{\textbf{Solving the Shortest Path Problem using Bellman's Equation}}: This problem involves using Bellman's equation and an iterative approximation method to find the shortest path in a graph, calculating shortest distances from each node to the destination node. <L[2024-11-14-NetAppls] 16><L[2024-11-14-NetAppls] 17><L[2024-11-14-NetAppls] 18><L[2024-11-14-NetAppls] 19>%
\item \href{https://www.math.purdue.edu/~yipn/421/2024-11-14-NetAppls.pdf#page=20}{\textbf{Solving the Shortest Path Problem using Dijkstra's Algorithm}}: This problem involves using Dijkstra's algorithm to find the shortest path in a graph, iteratively updating shortest distance estimates and tracking the path. <L[2024-11-14-NetAppls] 20><L[2024-11-14-NetAppls] 21>%
\end{itemize}%
\begin{itemize}[[leftmargin=*]]%
\item \href{https://www.math.purdue.edu/~yipn/421/2024-11-19-MaxFlowMinCut.pdf#page=1}{\textbf{Finding the Maximum Flow in a Network}}: This problem involves determining the largest possible flow that can be sent from a source node to a sink node in a network, subject to capacity constraints on the edges. <L[2024-11-19-MaxFlowMinCut] 1><L[2024-11-19-MaxFlowMinCut] 2><L[2024-11-19-MaxFlowMinCut] 3><L[2024-11-19-MaxFlowMinCut] 4><L[2024-11-19-MaxFlowMinCut] 5><L[2024-11-19-MaxFlowMinCut] 6><L[2024-11-19-MaxFlowMinCut] 7>%
\item \href{https://www.math.purdue.edu/~yipn/421/2024-11-19-MaxFlowMinCut.pdf#page=2}{\textbf{Defining and Calculating the Capacity of a Cut}}: This problem focuses on identifying a cut in a network and computing its capacity, which is the sum of the capacities of the edges crossing the cut from the source side to the sink side. <L[2024-11-19-MaxFlowMinCut] 2><L[2024-11-19-MaxFlowMinCut] 3>%
\item \href{https://www.math.purdue.edu/~yipn/421/2024-11-19-MaxFlowMinCut.pdf#page=8}{\textbf{Proving the Max{-}Flow Min{-}Cut Theorem}}: This problem involves demonstrating that the maximum flow in a network is equal to the minimum capacity among all possible cuts in the network. <L[2024-11-19-MaxFlowMinCut] 8><L[2024-11-19-MaxFlowMinCut] 9><L[2024-11-19-MaxFlowMinCut] 16>%
\item \href{https://www.math.purdue.edu/~yipn/421/2024-11-19-MaxFlowMinCut.pdf#page=10}{\textbf{Implementing the Ford{-}Fulkerson Algorithm}}: This problem involves applying the Ford-Fulkerson algorithm to find the maximum flow in a network by iteratively finding augmenting paths and increasing the flow along those paths. <L[2024-11-19-MaxFlowMinCut] 10><L[2024-11-19-MaxFlowMinCut] 11><L[2024-11-19-MaxFlowMinCut] 12><L[2024-11-19-MaxFlowMinCut] 13>%
\item \href{https://www.math.purdue.edu/~yipn/421/2024-11-19-MaxFlowMinCut.pdf#page=14}{\textbf{Formulating the Maximum Flow Problem as a Linear Program}}: This problem involves expressing the maximum flow problem as a linear program, defining the objective function and constraints in terms of flow variables and edge capacities. <L[2024-11-19-MaxFlowMinCut] 14><L[2024-11-19-MaxFlowMinCut] 15>%
\end{itemize}%
\begin{itemize}[[leftmargin=*]]%
\item \href{https://www.math.purdue.edu/~yipn/421/2024-11-21-EconNetSimplex.pdf#page=4}{\textbf{Finding a set of "fair prices" in a network flow problem}}: This problem involves determining a set of node prices ($y_i$) such that the price difference between two connected nodes equals the cost of transportation between them ($y_i + c_{ij} = y_j$) for all arcs in a spanning tree. <L[2024-11-21-EconNetSimplex] 4>%
\item \href{https://www.math.purdue.edu/~yipn/421/2024-11-21-EconNetSimplex.pdf#page=5}{\textbf{Identifying profitable arbitrage opportunities in a network flow problem}}: This problem focuses on finding arcs ($ij$) where the cost of buying at node $i$ and shipping to node $j$ is less than the selling price at node $j$ ($y_i + c_{ij} < y_j$), indicating a potential for profit improvement. <L[2024-11-21-EconNetSimplex] 5>%
\item \href{https://www.math.purdue.edu/~yipn/421/2024-11-21-EconNetSimplex.pdf#page=6}{\textbf{Adjusting shipments to maintain feasibility in a network flow problem}}: This problem involves modifying the flow along the arcs of a spanning tree to accommodate the introduction of a new arc, ensuring that the flow balance at each node is maintained. <L[2024-11-21-EconNetSimplex] 6>%
\item \href{https://www.math.purdue.edu/~yipn/421/2024-11-21-EconNetSimplex.pdf#page=3}{\textbf{Finding a feasible solution on a spanning tree}}: This problem involves finding an initial feasible solution for a network flow problem that is represented by a spanning tree, which is a subset of edges that connects all nodes without forming cycles. <L[2024-11-21-EconNetSimplex] 3>%
\end{itemize}%
\begin{itemize}[[leftmargin=*]]%
\item \href{https://www.math.purdue.edu/~yipn/421/2024-12-03-Duality.pdf#page=1}{\textbf{Formulating the Dual Problem}}: This problem involves constructing the dual linear programming problem from a given primal problem, correctly transposing coefficients and reversing inequality signs. <L[2024-12-03-Duality] 1><L[2024-12-03-Duality] 5><L[2024-12-03-Duality] 6><L[2024-12-03-Duality] 7><L[2024-12-03-Duality] 29><L[2024-12-03-Duality] 30>%
\item \href{https://www.math.purdue.edu/~yipn/421/2024-12-03-Duality.pdf#page=2}{\textbf{Applying Complementary Slackness}}: This problem focuses on using complementary slackness conditions to determine optimal solutions for both primal and dual problems. <L[2024-12-03-Duality] 2>%
\item \href{https://www.math.purdue.edu/~yipn/421/2024-12-03-Duality.pdf#page=3}{\textbf{Solving Linear Programs Graphically}}: This problem involves using a graphical method to solve linear programming problems with two variables by identifying the feasible region and the optimal solution point. <L[2024-12-03-Duality] 3><L[2024-12-03-Duality] 4><L[2024-12-03-Duality] 17><L[2024-12-03-Duality] 18><L[2024-12-03-Duality] 19>%
\item \href{https://www.math.purdue.edu/~yipn/421/2024-12-03-Duality.pdf#page=5}{\textbf{Formulating Primal and Dual Problems from a Given Problem Statement}}: This problem involves translating a word problem into its mathematical representation as both a primal and dual linear program. <L[2024-12-03-Duality] 5><L[2024-12-03-Duality] 6><L[2024-12-03-Duality] 7><L[2024-12-03-Duality] 29><L[2024-12-03-Duality] 30>%
\item \href{https://www.math.purdue.edu/~yipn/421/2024-12-03-Duality.pdf#page=8}{\textbf{Finding Optimal Solutions and Slack Variables}}: This problem involves finding the optimal solution for both the primal and dual problems and determining the values of the corresponding slack variables at optimality. <L[2024-12-03-Duality] 8>%
\item \href{https://www.math.purdue.edu/~yipn/421/2024-12-03-Duality.pdf#page=9}{\textbf{Geometrically Interpreting Dual Constraints}}: This problem involves interpreting the dual constraints geometrically as half-spaces in the primal variable space. <L[2024-12-03-Duality] 9><L[2024-12-03-Duality] 10>%
\item \href{https://www.math.purdue.edu/~yipn/421/2024-12-03-Duality.pdf#page=11}{\textbf{Performing Sensitivity Analysis on Objective Function Coefficients}}: This problem involves determining the range of changes in the objective function coefficients that will not alter the optimal solution. <L[2024-12-03-Duality] 11><L[2024-12-03-Duality] 12><L[2024-12-03-Duality] 13><L[2024-12-03-Duality] 14><L[2024-12-03-Duality] 15><L[2024-12-03-Duality] 20><L[2024-12-03-Duality] 21>%
\item \href{https://www.math.purdue.edu/~yipn/421/2024-12-03-Duality.pdf#page=11}{\textbf{Performing Sensitivity Analysis on Constraint Values}}: This problem involves determining the range of changes in the constraint values that will not alter the optimal solution. <L[2024-12-03-Duality] 11><L[2024-12-03-Duality] 12><L[2024-12-03-Duality] 13><L[2024-12-03-Duality] 14><L[2024-12-03-Duality] 15><L[2024-12-03-Duality] 20><L[2024-12-03-Duality] 21><L[2024-12-03-Duality] 26><L[2024-12-03-Duality] 27>%
\item \href{https://www.math.purdue.edu/~yipn/421/2024-12-03-Duality.pdf#page=15}{\textbf{Sensitivity Analysis using Shadow Prices}}: This problem involves using shadow prices (optimal dual variables) to determine the impact of changes in constraint values on the optimal objective function value. <L[2024-12-03-Duality] 15><L[2024-12-03-Duality] 20><L[2024-12-03-Duality] 21><L[2024-12-03-Duality] 25><L[2024-12-03-Duality] 26><L[2024-12-03-Duality] 27>%
\item \href{https://www.math.purdue.edu/~yipn/421/2024-12-03-Duality.pdf#page=16}{\textbf{Solving a Linear Program and Performing Sensitivity Analysis}}: This problem involves solving a linear program and then determining the range of changes in both the objective function coefficients and constraint values that maintain the optimality of the initial solution. <L[2024-12-03-Duality] 16>%
\item \href{https://www.math.purdue.edu/~yipn/421/2024-12-03-Duality.pdf#page=22}{\textbf{Setting up a Linear Program in Matrix Form}}: This problem involves representing a linear program in matrix form, identifying the basis and non-basic variables. <L[2024-12-03-Duality] 22><L[2024-12-03-Duality] 23><L[2024-12-03-Duality] 24><L[2024-12-03-Duality] 25>%
\item \href{https://www.math.purdue.edu/~yipn/421/2024-12-03-Duality.pdf#page=25}{\textbf{Determining Optimal Solution Value and Changes in Optimal Solution Value}}: This problem involves calculating the optimal solution value and determining how changes in the constraint values affect this value. <L[2024-12-03-Duality] 25>%
\item \href{https://www.math.purdue.edu/~yipn/421/2024-12-03-Duality.pdf#page=31}{\textbf{Analyzing Complementary Slackness Conditions}}: This problem involves analyzing the complementary slackness conditions to determine the values of slack variables and their implications for the optimal solution. <L[2024-12-03-Duality] 31><L[2024-12-03-Duality] 32><L[2024-12-03-Duality] 34><L[2024-12-03-Duality] 35>%
\item \href{https://www.math.purdue.edu/~yipn/421/2024-12-03-Duality.pdf#page=33}{\textbf{Formulating and Solving a Dual Problem in Sensitivity Analysis}}: This problem involves formulating and solving a dual problem to perform sensitivity analysis on a linear program. <L[2024-12-03-Duality] 33><L[2024-12-03-Duality] 36><L[2024-12-03-Duality] 37>%
\end{itemize}%
\newpage

%
\section{Algorithm Solutions}%
\label{sec:AlgorithmSolutions}%
\begin{itemize}[[leftmargin=*]]%
\item \href{https://www.math.purdue.edu/~yipn/421/2024-08-27-ExSimplex.pdf#page=6}{\textbf{Simplex Method}}: Iteratively move from one vertex of the feasible region to another by setting non-basic variables to zero and choosing a new set of (non)basic variables to improve the objective function, until the optimal solution is found.<L[2024-08-27-ExSimplex] 6><L[2024-08-27-ExSimplex] 7><L[2024-08-27-ExSimplex] 8><L[2024-08-27-ExSimplex] 9><L[2024-08-27-ExSimplex] 10><L[2024-08-27-ExSimplex] 11><L[2024-08-27-ExSimplex] 12><L[2024-08-27-ExSimplex] 13><L[2024-08-27-ExSimplex] 14><L[2024-08-27-ExSimplex] 15><L[2024-08-27-ExSimplex] 16><L[2024-08-27-ExSimplex] 17><L[2024-08-27-ExSimplex] 18><L[2024-08-27-ExSimplex] 19><L[2024-08-27-ExSimplex] 20><L[2024-08-27-ExSimplex] 21><L[2024-08-27-ExSimplex] 22><L[2024-08-27-ExSimplex] 23><L[2024-08-27-ExSimplex] 24><L[2024-08-27-ExSimplex] 25>%
\item \href{https://www.math.purdue.edu/~yipn/421/2024-08-27-ExSimplex.pdf#page=32}{\textbf{Auxiliary Problem}}: Introduce a new variable  $x_0$ and modify constraints to create a new problem whose feasible region includes the origin. Solve this problem to find a feasible solution for the original problem, then proceed with the Simplex method.<L[2024-08-27-ExSimplex] 32><L[2024-08-27-ExSimplex] 33><L[2024-08-27-ExSimplex] 34><L[2024-08-27-ExSimplex] 35><L[2024-08-27-ExSimplex] 36><L[2024-08-27-ExSimplex] 37><L[2024-08-27-ExSimplex] 38><L[2024-08-27-ExSimplex] 39><L[2024-08-27-ExSimplex] 40><L[2024-08-27-ExSimplex] 41><L[2024-08-27-ExSimplex] 42><L[2024-08-27-ExSimplex] 43>%
\end{itemize}%
\begin{itemize}[[leftmargin=*]]%
\item \href{https://www.math.purdue.edu/~yipn/421/2024-09-05-SimplexEfficiency.pdf#page=1}{\textbf{Simplex Method Efficiency}}: The number of iterations in the Simplex method is empirically bounded by $\frac{3m}{2}$ for $m < 50$ and $m+n < 200$, where m is the number of constraints and n is the number of variables. <L[2024-09-05-SimplexEfficiency] 1>%
\item \href{https://www.math.purdue.edu/~yipn/421/2024-09-05-SimplexEfficiency.pdf#page=6}{\textbf{Largest{-}Increase Rule}}: A pivot rule for the Simplex method that selects the entering variable that yields the largest increase in the objective function value. <L[2024-09-05-SimplexEfficiency] 6><L[2024-09-05-SimplexEfficiency] 7><L[2024-09-05-SimplexEfficiency] 8>%
\item \href{https://www.math.purdue.edu/~yipn/421/2024-09-05-SimplexEfficiency.pdf#page=10}{\textbf{Randomized Pivot Rule}}: A pivot rule that randomly permutes the order of variables at the start, then uses Bland's rule for entering variables and a lexicographic method for leaving variables. The expected number of pivot steps is bounded by $e C \sqrt{n} \ln n$. <L[2024-09-05-SimplexEfficiency] 10>%
\item \href{https://www.math.purdue.edu/~yipn/421/2024-09-05-SimplexEfficiency.pdf#page=11}{\textbf{Smoothed Complexity}}: The analysis of the Simplex method's performance when the input linear program's coefficients are perturbed by small random amounts. With high probability, the number of steps is polynomial in the problem size. <L[2024-09-05-SimplexEfficiency] 11>%
\item \href{https://www.math.purdue.edu/~yipn/421/2024-09-05-SimplexEfficiency.pdf#page=12}{\textbf{Interior Point Methods}}: Polynomial-time algorithms for linear programming that traverse the interior of the feasible region, unlike the Simplex method which follows the boundary. <L[2024-09-05-SimplexEfficiency] 12>%
\end{itemize}%
\begin{itemize}[[leftmargin=*]]%
\item \href{https://www.math.purdue.edu/~yipn/421/2024-09-12-Duality.pdf#page=2}{\textbf{Weak Duality}}: $\sum_{j} c_j x_j \le \sum_{i} b_i y_i$<L[2024-09-12-Duality] 2>%
\item \textbf{Simplex Method}: Iteratively improves a feasible solution by exchanging basic and non-basic variables until optimality is reached (Slides 6, 7, 8, 9, 10). The negative transpose property is preserved in each step.%
\item \href{https://www.math.purdue.edu/~yipn/421/2024-09-12-Duality.pdf#page=11}{\textbf{Strong Duality}}: If (P) has an optimal solution ($x_j^$), then (D) also has an optimal solution ($y_i^$) and $\sum_{j} c_j x_j^ = \sum_{i} b_i y_i^$<L[2024-09-12-Duality] 11><L[2024-09-12-Duality] 17>%
\item \href{https://www.math.purdue.edu/~yipn/421/2024-09-12-Duality.pdf#page=21}{\textbf{Complementary Slackness}}: $x_j z_j = 0$ for $j = 1, 2, \dots, n$ and $w_i y_i = 0$ for $i = 1, 2, \dots, m$<L[2024-09-12-Duality] 21>%
\end{itemize}%
\begin{itemize}[[leftmargin=*]]%
\item \href{https://www.math.purdue.edu/~yipn/421/2024-09-17-DualityGeneral.pdf#page=7}{\textbf{Lagrangian Function}}: $L(x,\lambda,\mu) = f_0(x) - \sum_{i=1}^m \lambda_i f_i(x) - \sum_{j=1}^n \mu_j g_j(x)$<L[2024-09-17-DualityGeneral] 7>%
\item \href{https://www.math.purdue.edu/~yipn/421/2024-09-17-DualityGeneral.pdf#page=5}{\textbf{Weak Duality}}: $\max_{x \in \mathcal{G}} \xi(x) \le \min_{\lambda \ge 0, \mu} \xi(\lambda,\mu)$<L[2024-09-17-DualityGeneral] 5><L[2024-09-17-DualityGeneral] 9><L[2024-09-17-DualityGeneral] 10><L[2024-09-17-DualityGeneral] 18>%
\item \href{https://www.math.purdue.edu/~yipn/421/2024-09-17-DualityGeneral.pdf#page=17}{\textbf{Strong Duality}}: If $x^$ is optimal for (P) and $\lambda^$ is optimal for (D), then $\xi(x^) = \xi(\lambda^)$<L[2024-09-17-DualityGeneral] 17><L[2024-09-17-DualityGeneral] 21><L[2024-09-17-DualityGeneral] 22>%
\item \href{https://www.math.purdue.edu/~yipn/421/2024-09-17-DualityGeneral.pdf#page=22}{\textbf{Complementary Slackness}}: If $x^$ is optimal for (P) and $\lambda^$ is optimal for (D), then $x^_i \lambda^_i = 0$ for all $i$<L[2024-09-17-DualityGeneral] 22>%
\end{itemize}%
\begin{itemize}[[leftmargin=*]]%
\item \href{https://www.math.purdue.edu/~yipn/421/2024-09-19-DualityExamples.pdf#page=5}{\textbf{Simplex Algorithm}}: This algorithm iteratively improves a feasible solution by moving from one vertex of the feasible region to an adjacent vertex with a better objective function value until an optimal solution is found. <L[2024-09-19-DualityExamples] 5><L[2024-09-19-DualityExamples] 8>%
\item \href{https://www.math.purdue.edu/~yipn/421/2024-09-19-DualityExamples.pdf#page=8}{\textbf{Dual Simplex Algorithm}}: This algorithm starts with a dual feasible solution and iteratively improves it by moving to adjacent vertices with better objective function values until an optimal solution is found. <L[2024-09-19-DualityExamples] 8>%
\end{itemize}%
\begin{itemize}[[leftmargin=*]]%
\item \href{https://www.math.purdue.edu/~yipn/421/2024-09-24-SimplexMatrix.pdf#page=8}{\textbf{Simplex Method}}: Iteratively improves a feasible solution by exchanging basic and non-basic variables until an optimal solution is found. The process involves solving for basic variables ($\mathbf{x}_B = B^{-1}\mathbf{b} - B^{-1}N\mathbf{x}_N$) and substituting into the objective function ($\mathcal{Z} = \mathbf{c}_B^T B^{-1} \mathbf{b} + (\mathbf{c}_N^T - \mathbf{c}_B^T B^{-1} N) \mathbf{x}_N$) to identify variables to enter and leave the basis. <L[2024-09-24-SimplexMatrix] 8><L[2024-09-24-SimplexMatrix] 9><L[2024-09-24-SimplexMatrix] 10><L[2024-09-24-SimplexMatrix] 11><L[2024-09-24-SimplexMatrix] 12><L[2024-09-24-SimplexMatrix] 13><L[2024-09-24-SimplexMatrix] 17><L[2024-09-24-SimplexMatrix] 18><L[2024-09-24-SimplexMatrix] 19><L[2024-09-24-SimplexMatrix] 20>%
\item \href{https://www.math.purdue.edu/~yipn/421/2024-09-24-SimplexMatrix.pdf#page=4}{\textbf{Basic and Non{-}basic Variables}}: Partitions the variables into basic ($\mathbf{x}_B$) and non-basic ($\mathbf{x}_N$) sets. Basic variables are linearly independent and determine the current feasible solution. Non-basic variables are set to zero. <L[2024-09-24-SimplexMatrix] 4><L[2024-09-24-SimplexMatrix] 5><L[2024-09-24-SimplexMatrix] 6><L[2024-09-24-SimplexMatrix] 7>%
\item \href{https://www.math.purdue.edu/~yipn/421/2024-09-24-SimplexMatrix.pdf#page=14}{\textbf{Complementary Slackness}}: Relates primal and dual variables. If a primal variable is non-zero, its corresponding dual variable is zero, and vice versa. <L[2024-09-24-SimplexMatrix] 14>%
\end{itemize}%
\begin{itemize}[[leftmargin=*]]%
\item \href{https://www.math.purdue.edu/~yipn/421/2024-09-26-Sensitivity.pdf#page=2}{\textbf{Simplex Method}}: This algorithm iteratively improves a feasible solution to a linear program until an optimal solution is found. It moves from one vertex of the feasible region to an adjacent vertex with a better objective function value until no further improvement is possible.<L[2024-09-26-Sensitivity] 2><L[2024-09-26-Sensitivity] 3><L[2024-09-26-Sensitivity] 7><L[2024-09-26-Sensitivity] 11>%
\item \href{https://www.math.purdue.edu/~yipn/421/2024-09-26-Sensitivity.pdf#page=4}{\textbf{Sensitivity Analysis}}: This analysis determines the range of changes in the objective function coefficients or constraint values that will not change the optimal solution. It uses the optimal simplex tableau to calculate the allowable changes.<L[2024-09-26-Sensitivity] 4><L[2024-09-26-Sensitivity] 5><L[2024-09-26-Sensitivity] 6><L[2024-09-26-Sensitivity] 7><L[2024-09-26-Sensitivity] 8><L[2024-09-26-Sensitivity] 9><L[2024-09-26-Sensitivity] 10><L[2024-09-26-Sensitivity] 11>%
\item \href{https://www.math.purdue.edu/~yipn/421/2024-09-26-Sensitivity.pdf#page=4}{\textbf{Parametric Analysis}}: This analysis investigates how the optimal solution changes as a parameter in the objective function or constraints varies continuously. It involves solving a sequence of linear programs with incrementally changing parameters.<L[2024-09-26-Sensitivity] 4><L[2024-09-26-Sensitivity] 8>%
\end{itemize}%
\begin{itemize}[[leftmargin=*]]%
\item \href{https://www.math.purdue.edu/~yipn/421/2024-10-01-DualGeneralLP.pdf#page=1}{\textbf{Dual of General LP Problem}}: Given a primal problem $\max \vec{s} = \vec{c}\vec{x}$ subject to $\vec{A}\vec{x} \le \vec{b}$ and $\vec{x} \ge \vec{0}$, its dual is $\min \vec{s} = \vec{b}\vec{y}$ subject to $\vec{A}^T\vec{y} \ge \vec{c}$ and $\vec{y} \ge \vec{0}$.<L[2024-10-01-DualGeneralLP] 1>%
\item \href{https://www.math.purdue.edu/~yipn/421/2024-10-01-DualGeneralLP.pdf#page=2}{\textbf{Reformulation of Primal Problem with Unrestricted Variables}}: To solve a primal problem $\max \vec{c}^T\vec{x}$ subject to $\vec{A}\vec{x} \le \vec{b}$ and $\vec{x} \le \vec{u}$ (where $\vec{x}$ is not necessarily non-negative), rewrite it as $\max \vec{c}^T(\vec{x}^+ - \vec{x}^-)$ subject to $\vec{A}(\vec{x}^+ - \vec{x}^-) \le \vec{b}$, $\vec{x}^+, \vec{x}^- \ge \vec{0}$, and $\vec{x} = \vec{x}^+ - \vec{x}^-$.<L[2024-10-01-DualGeneralLP] 2>%
\item \href{https://www.math.purdue.edu/~yipn/421/2024-10-01-DualGeneralLP.pdf#page=3}{\textbf{Dual of General LP with Upper Bounds}}: The dual of $\max \vec{c}^T\vec{x}$ subject to $\vec{A}\vec{x} \le \vec{b}$ and $\vec{x} \le \vec{u}$ is $\min \vec{b}^T\vec{y} + \vec{u}^T\vec{z}$ subject to $\vec{A}^T\vec{y} + \vec{z} \ge \vec{c}$ and $\vec{y}, \vec{z} \ge \vec{0}$.<L[2024-10-01-DualGeneralLP] 3>%
\item \href{https://www.math.purdue.edu/~yipn/421/2024-10-01-DualGeneralLP.pdf#page=4}{\textbf{Lagrangian Function for LP with Upper Bounds}}: The Lagrangian function for the primal problem $\max \vec{c}^T\vec{x}$ subject to $\vec{A}\vec{x} \le \vec{b}$ and $\vec{x} \le \vec{u}$ is $\mathcal{L}(\vec{x}, \vec{y}, \vec{z}) = \vec{c}^T\vec{x} - \vec{y}^T(\vec{A}\vec{x} - \vec{b}) - \vec{z}^T(\vec{x} - \vec{u})$, where $\vec{y} \ge \vec{0}$ and $\vec{z} \ge \vec{0}$.<L[2024-10-01-DualGeneralLP] 4>%
\item \href{https://www.math.purdue.edu/~yipn/421/2024-10-01-DualGeneralLP.pdf#page=6}{\textbf{Derivation of Dual Problem from Lagrangian}}: Minimizing the Lagrangian function $\mathcal{L}(\vec{x}, \vec{y}, \vec{z})$ with respect to $\vec{y}$ and $\vec{z}$ leads to the dual problem's constraints $\vec{A}^T\vec{y} + \vec{z} = \vec{c}$ and objective function $\min_{\vec{y}, \vec{z} \ge \vec{0}} \vec{y}^T\vec{b} + \vec{z}^T\vec{u}$.<L[2024-10-01-DualGeneralLP] 6><L[2024-10-01-DualGeneralLP] 7>%
\end{itemize}%
\begin{itemize}[[leftmargin=*]]%
\item \href{https://www.math.purdue.edu/~yipn/421/2024-10-01-SensitivityDuality.pdf#page=4}{\textbf{Simplex Method}}: This algorithm iteratively improves a feasible solution to a linear program by moving from one vertex of the feasible region to an adjacent vertex with a better objective function value until an optimal solution is found.<L[2024-10-01-SensitivityDuality] 4><L[2024-10-01-SensitivityDuality] 5><L[2024-10-01-SensitivityDuality] 6>%
\item \href{https://www.math.purdue.edu/~yipn/421/2024-10-01-SensitivityDuality.pdf#page=3}{\textbf{Complementary Slackness}}: If $x_i^ > 0$, then the $i$-th constraint in the dual problem is binding ($A^T y = c$). If the $i$-th constraint in the primal problem is not binding ($Ax < b$), then $y_i^ = 0$.<L[2024-10-01-SensitivityDuality] 3>%
\item \href{https://www.math.purdue.edu/~yipn/421/2024-10-01-SensitivityDuality.pdf#page=2}{\textbf{Strong Duality}}: The optimal objective function values of the primal and dual problems are equal.<L[2024-10-01-SensitivityDuality] 2>%
\item \href{https://www.math.purdue.edu/~yipn/421/2024-10-01-SensitivityDuality.pdf#page=1}{\textbf{Sensitivity Analysis}}: This technique examines how changes in the parameters of a linear program (e.g., objective function coefficients, constraint coefficients, or right-hand side values) affect the optimal solution.<L[2024-10-01-SensitivityDuality] 1><L[2024-10-01-SensitivityDuality] 7><L[2024-10-01-SensitivityDuality] 8><L[2024-10-01-SensitivityDuality] 9><L[2024-10-01-SensitivityDuality] 10>%
\item \href{https://www.math.purdue.edu/~yipn/421/2024-10-01-SensitivityDuality.pdf#page=10}{\textbf{Shadow Price}}: The rate of change in the optimal objective function value for a unit change in the right-hand side of a constraint. It is given by the optimal dual variable corresponding to that constraint.<L[2024-10-01-SensitivityDuality] 10>%
\end{itemize}%
\begin{itemize}[[leftmargin=*]]%
\item \href{https://www.math.purdue.edu/~yipn/421/2024-10-03-DualSimplexMatrix.pdf#page=1}{\textbf{Strong Duality}}: This theorem states that the optimal objective function values of the primal and dual problems are equal.<L[2024-10-03-DualSimplexMatrix] 1>%
\item \href{https://www.math.purdue.edu/~yipn/421/2024-10-03-DualSimplexMatrix.pdf#page=2}{\textbf{Dual Simplex Method}}: This method solves linear programming problems by iteratively improving the dual feasible solution until an optimal solution is found.<L[2024-10-03-DualSimplexMatrix] 2><L[2024-10-03-DualSimplexMatrix] 3><L[2024-10-03-DualSimplexMatrix] 4>%
\item \href{https://www.math.purdue.edu/~yipn/421/2024-10-03-DualSimplexMatrix.pdf#page=5}{\textbf{Dual Based Phase I Algorithm}}: This algorithm finds a feasible solution for the dual problem, which can then be used to find a feasible solution for the primal problem.<L[2024-10-03-DualSimplexMatrix] 5><L[2024-10-03-DualSimplexMatrix] 6>%
\item \href{https://www.math.purdue.edu/~yipn/421/2024-10-03-DualSimplexMatrix.pdf#page=7}{\textbf{Simplex Method in Matrix Form}}: This method represents the simplex algorithm using matrices, allowing for efficient computation and analysis.<L[2024-10-03-DualSimplexMatrix] 7><L[2024-10-03-DualSimplexMatrix] 8><L[2024-10-03-DualSimplexMatrix] 9><L[2024-10-03-DualSimplexMatrix] 10><L[2024-10-03-DualSimplexMatrix] 11><L[2024-10-03-DualSimplexMatrix] 12><L[2024-10-03-DualSimplexMatrix] 13><L[2024-10-03-DualSimplexMatrix] 14><L[2024-10-03-DualSimplexMatrix] 15>%
\end{itemize}%
\begin{itemize}[[leftmargin=*]]%
\item \href{https://www.math.purdue.edu/~yipn/421/2024-10-22-Convexity.pdf#page=1}{\textbf{Convex Set}}: $S \subseteq \mathbb{R}^m$ is convex if for any $z_1, z_2 \in S \implies \lambda z_1 + (1-\lambda) z_2 \in S$, $0 \le \lambda \le 1$.<L[2024-10-22-Convexity] 1>%
\item \href{https://www.math.purdue.edu/~yipn/421/2024-10-22-Convexity.pdf#page=1}{\textbf{Convex Combination}}: $\sum_{i=1}^n \lambda_i z_i$, $\lambda_i \ge 0$, $\sum_{i=1}^n \lambda_i = 1$.<L[2024-10-22-Convexity] 1>%
\item \href{https://www.math.purdue.edu/~yipn/421/2024-10-22-Convexity.pdf#page=2}{\textbf{Theorem 10.1}}: $S$ is convex $\iff$ it contains all convex combinations of points in $S$.<L[2024-10-22-Convexity] 2><L[2024-10-22-Convexity] 3><L[2024-10-22-Convexity] 4>%
\item \href{https://www.math.purdue.edu/~yipn/421/2024-10-22-Convexity.pdf#page=5}{\textbf{Convex Hull}}: $\text{Conv}(S) = \bigcap \{C \text{ convex } | S \subseteq C\}$ (= Smallest convex set containing $S$).<L[2024-10-22-Convexity] 5>%
\item \href{https://www.math.purdue.edu/~yipn/421/2024-10-22-Convexity.pdf#page=5}{\textbf{Theorem 10.2}}: $S \subseteq \mathbb{R}^m$, $\text{Conv}(S) = \{\text{collection of all convex combinations of finitely many points from } S\}$.<L[2024-10-22-Convexity] 5><L[2024-10-22-Convexity] 6><L[2024-10-22-Convexity] 7><L[2024-10-22-Convexity] 8>%
\item \href{https://www.math.purdue.edu/~yipn/421/2024-10-22-Convexity.pdf#page=8}{\textbf{Caratheodory's Theorem}}: If $z \in \text{Conv}(S)$, then $z$ can be written as a convex combination of at most $m+1$ points from $S$.<L[2024-10-22-Convexity] 8><L[2024-10-22-Convexity] 9>%
\item \href{https://www.math.purdue.edu/~yipn/421/2024-10-22-Convexity.pdf#page=10}{\textbf{Separation Theorem}}: Let $P$ and $\bar{P}$ be two disjoint nonempty polyhedra in $\mathbb{R}^n$. Then there exist disjoint half-spaces $H$ and $\bar{H}$ such that $P \subset H$ and $\bar{P} \subset \bar{H}$.<L[2024-10-22-Convexity] 10><L[2024-10-22-Convexity] 11><L[2024-10-22-Convexity] 12><L[2024-10-22-Convexity] 13><L[2024-10-22-Convexity] 14><L[2024-10-22-Convexity] 15>%
\item \href{https://www.math.purdue.edu/~yipn/421/2024-10-22-Convexity.pdf#page=12}{\textbf{Farkas' Lemma}}: The system $Ax \le b$ has no solution if and only if there exists a vector $y \ge 0$ such that $y^T A = 0$ and $y^T b < 0$.<L[2024-10-22-Convexity] 12><L[2024-10-22-Convexity] 13><L[2024-10-22-Convexity] 14><L[2024-10-22-Convexity] 15>%
\end{itemize}%
\begin{itemize}[[leftmargin=*]]%
\item \href{https://www.math.purdue.edu/~yipn/421/2024-10-22-Regression.pdf#page=4}{\textbf{\$L\^{}2\$ Regression (Least Square)}}: $\min_{(x_1, \dots, x_n)} \sum_{i=1}^m \left| (b_i - \sum_{j=1}^n a_{ij} x_j) \right|^2$ or $\min \left\| b - AX \right\|_2^2$, solved by the normal equation $A^T A \hat{X} = A^T b$, yielding $\hat{X} = (A^T A)^{-1} A^T b$ if $(A^T A)^{-1}$ exists.<L[2024-10-22-Regression] 4>%
\item \href{https://www.math.purdue.edu/~yipn/421/2024-10-22-Regression.pdf#page=5}{\textbf{\$L\^{}1\$ Regression}}: $\min_{x_1, \dots, x_n} \sum_{i=1}^m |b_i - \sum_{j=1}^n a_{ij} x_j|$ or $\min_X \|b - AX\|_1$, solved by introducing slack variables $\hat{\varepsilon}_i$ such that $-\hat{\varepsilon}_i \le b_i - \sum_{j=1}^n a_{ij} x_j \le \hat{\varepsilon}_i$ and minimizing $\sum_{i=1}^m \hat{\varepsilon}_i$.<L[2024-10-22-Regression] 5>%
\item \href{https://www.math.purdue.edu/~yipn/421/2024-10-22-Regression.pdf#page=6}{\textbf{\$L\^{}\textbackslash{}infty\$ Regression}}: $\min_{x_1, \dots, x_n} \max_{i} |b_i - \sum_{j=1}^n a_{ij} x_j|$ or $\min_X \|b - AX\|_\infty$, solved by introducing a slack variable $\delta$ such that $-\delta \le b_i - \sum_{j=1}^n a_{ij} x_j \le \delta$ for all $i$ and minimizing $\delta$.<L[2024-10-22-Regression] 6>%
\item \href{https://www.math.purdue.edu/~yipn/421/2024-10-22-Regression.pdf#page=11}{\textbf{Lagrangian Formulation of SVM}}: $\min_{\alpha} \frac{1}{2} \sum_{i=1}^N \sum_{j=1}^N \alpha_i \alpha_j y_i y_j x_i^T x_j - \sum_{i=1}^N \alpha_i$ subject to $\sum_{i=1}^N \alpha_i y_i = 0$ and $0 \le \alpha_i \le C$ for all $i$.<L[2024-10-22-Regression] 11>%
\item \href{https://www.math.purdue.edu/~yipn/421/2024-10-22-Regression.pdf#page=13}{\textbf{SVM Optimization as a QP Problem (version 1)}}: Maximize $\delta$ subject to $y_i (a_1 x_i^1 + a_2 x_i^2 - b) \ge \delta$ for all $i$.<L[2024-10-22-Regression] 13>%
\item \href{https://www.math.purdue.edu/~yipn/421/2024-10-22-Regression.pdf#page=14}{\textbf{SVM Optimization as a QP Problem (version 2)}}: Maximize $\delta$ subject to $y_i (a_1 x_i^1 + a_2 x_i^2 - b) \ge \delta$ for all $i$ and $-1 \le a_1, a_2 \le 1$.<L[2024-10-22-Regression] 14>%
\item \href{https://www.math.purdue.edu/~yipn/421/2024-10-22-Regression.pdf#page=19}{\textbf{SVM Optimization as a QP Problem (version 3)}}: Maximize $d$ subject to $|a_1 s_i + a_2 s_i^2 - b'| \le d \sqrt{a_1^2 + a_2^2}$ for all $i$ and $a_1 s_i + a_2 s_i^2 - b' \le 0$ for all $i$.<L[2024-10-22-Regression] 19>%
\item \href{https://www.math.purdue.edu/~yipn/421/2024-10-22-Regression.pdf#page=20}{\textbf{SVM Optimization as a QP Problem (version 4)}}: Maximize $d$ subject to $d \le \frac{|a_1 s_i + a_2 s_i^2 - b'|}{\sqrt{a_1^2 + a_2^2}}$ for all $i$ and $a_1 s_i + a_2 s_i^2 - b' \le 0$ for all $i$.<L[2024-10-22-Regression] 20>%
\item \href{https://www.math.purdue.edu/~yipn/421/2024-10-22-Regression.pdf#page=22}{\textbf{Smallest Ball Problem}}: Given points $p_1, \dots, p_n \in \mathbb{R}^d$, find a ball of the smallest radius that contains all the points.<L[2024-10-22-Regression] 22>%
\end{itemize}%
\begin{itemize}[[leftmargin=*]]%
\item \href{https://www.math.purdue.edu/~yipn/421/2024-10-28-ApplsLinProg.pdf#page=5}{\textbf{Ice Cream Production Planning}}: $\min 50 \sum_{i=1}^{12} (Y_i + Z_i) + 20 \sum_{i=1}^{12} S_i$  s.t. $X_i + S_{i-1} - S_i = d_i$, $X_i - X_{i-1} = Y_i - Z_i$, $X_0 = 0$, $S_0 = 0$, $S_{12} = 0$, $X_i, Y_i, Z_i, S_i \ge 0$ <L[2024-10-28-ApplsLinProg] 5>%
\item \href{https://www.math.purdue.edu/~yipn/421/2024-10-28-ApplsLinProg.pdf#page=8}{\textbf{Workforce Planning}}: $\min 25 \sum_{i=1}^{12} Y_i + (17.5 \times 12) X$ s.t. $a_i - X \le Y_i$, $0 \le Y_i$, $X \ge 0$ <L[2024-10-28-ApplsLinProg] 8>%
\item \href{https://www.math.purdue.edu/~yipn/421/2024-10-28-ApplsLinProg.pdf#page=16}{\textbf{Markowitz Portfolio Selection Model}}: $\max \mu \sum_{j} x_j \overline{r_j} - \frac{\epsilon}{2} \sum_{i} \sum_{j} x_i x_j \sigma_{ij}$ s.t. $\sum_{j} x_j = 1$, $x_j \ge 0$ <L[2024-10-28-ApplsLinProg] 16>%
\item \href{https://www.math.purdue.edu/~yipn/421/2024-10-28-ApplsLinProg.pdf#page=18}{\textbf{Smallest Enclosing Ball Problem}}: minimize $x^T Q^T Q x - \sum_{j=1}^n x_j p_j^T p_j$ subject to $\sum_{j=1}^n x_j = 1$, $x_j \ge 0$ <L[2024-10-28-ApplsLinProg] 18>%
\item \href{https://www.math.purdue.edu/~yipn/421/2024-10-28-ApplsLinProg.pdf#page=21}{\textbf{Maximum Weight Matching}}: $\max \sum_{e \in E} w_e x_e$ s.t. $\sum_{e \in \delta(v)} x_e = 1$ for all vertices $v$, $x_e \in \{0, 1\}$ for all edges $e$ <L[2024-10-28-ApplsLinProg] 21>%
\item \href{https://www.math.purdue.edu/~yipn/421/2024-10-28-ApplsLinProg.pdf#page=23}{\textbf{Machine Scheduling}}: $\min \pi$ s.t. $\sum_{i \in M} x_{ij} = 1$ for $j \in J$, $\sum_{j \in J} d_{ij} x_{ij} \le \pi$ for $i \in M$, $x_{ij} \in \{0, 1\}$ for all $i, j$ <L[2024-10-28-ApplsLinProg] 23>%
\item \href{https://www.math.purdue.edu/~yipn/421/2024-10-28-ApplsLinProg.pdf#page=24}{\textbf{Knapsack Problem}}: maximize $\sum_{j=1}^n c_j x_j$ subject to $\sum_{j=1}^n a_j x_j \le b$, $x_j \in \{0, 1\}$, $j = 1, 2, \dots, n$ <L[2024-10-28-ApplsLinProg] 24>%
\item \href{https://www.math.purdue.edu/~yipn/421/2024-10-28-ApplsLinProg.pdf#page=29}{\textbf{Cutting Stock Problem}}: $\min x_1 + x_2 + \dots + x_{12}$ subject to  constraints based on cutting patterns $P_i$ and demand for each roll width. <L[2024-10-28-ApplsLinProg] 29>%
\item \href{https://www.math.purdue.edu/~yipn/421/2024-10-28-ApplsLinProg.pdf#page=34}{\textbf{Sparse Solutions of Linear System}}: $\min \|x\|_1$ s.t. $Ax = b$ <L[2024-10-28-ApplsLinProg] 34>%
\end{itemize}%
\begin{itemize}[[leftmargin=*]]%
\item \href{https://www.math.purdue.edu/~yipn/421/2024-10-28-Farkas.pdf#page=1}{\textbf{Farkas Lemma}}: The system $Ax \le b$ has no solutions if and only if there is a $y$ such that $A^T y = 0$, $y \ge 0$, and $b^T y < 0$.<L[2024-10-28-Farkas] 1>%
\item \href{https://www.math.purdue.edu/~yipn/421/2024-10-28-Farkas.pdf#page=4}{\textbf{Farkas Lemma (Three Variants)}}: (i) $Ax = b$ has a nonnegative solution $x \ge 0$ iff every $y$ with $y^T A \ge 0^T$ satisfies $y^T b \ge 0$. (ii) $Ax \le b$ has a nonnegative solution $x \ge 0$ iff every $y$ with $y^T A \ge 0^T$ satisfies $y^T b \ge 0$. (iii) $Ax \le b$ has a solution $x$ iff every $y$ with $y^T A = 0^T$ satisfies $y^T b \ge 0$.<L[2024-10-28-Farkas] 4><L[2024-10-28-Farkas] 5><L[2024-10-28-Farkas] 6><L[2024-10-28-Farkas] 7><L[2024-10-28-Farkas] 8><L[2024-10-28-Farkas] 9><L[2024-10-28-Farkas] 10><L[2024-10-28-Farkas] 11>%
\end{itemize}%
\begin{itemize}[[leftmargin=*]]%
\item \href{https://www.math.purdue.edu/~yipn/421/2024-10-31-IntLinProg.pdf#page=5}{\textbf{Branch and Bound}}: This algorithm solves integer programming problems by recursively partitioning the feasible region into subproblems, solving the relaxed linear program for each subproblem, and using the objective function values to prune branches that cannot lead to a better integer solution. A depth-first search strategy is often employed to find feasible integer solutions early, enabling pruning and improving efficiency. <L[2024-10-31-IntLinProg] 5><L[2024-10-31-IntLinProg] 6><L[2024-10-31-IntLinProg] 7><L[2024-10-31-IntLinProg] 8><L[2024-10-31-IntLinProg] 9><L[2024-10-31-IntLinProg] 10><L[2024-10-31-IntLinProg] 11><L[2024-10-31-IntLinProg] 12><L[2024-10-31-IntLinProg] 13><L[2024-10-31-IntLinProg] 14>%
\item \href{https://www.math.purdue.edu/~yipn/421/2024-10-31-IntLinProg.pdf#page=15}{\textbf{Gomory Cuts}}: This algorithm adds constraints (cuts) to the relaxed linear program to eliminate non-integer solutions. The cuts are derived from the fractional parts of the coefficients in the optimal dictionary of the relaxed problem. The process is iterative, adding cuts until an integer solution is found. <L[2024-10-31-IntLinProg] 15><L[2024-10-31-IntLinProg] 16><L[2024-10-31-IntLinProg] 17><L[2024-10-31-IntLinProg] 18><L[2024-10-31-IntLinProg] 19><L[2024-10-31-IntLinProg] 20><L[2024-10-31-IntLinProg] 21><L[2024-10-31-IntLinProg] 22><L[2024-10-31-IntLinProg] 23>%
\end{itemize}%
\begin{itemize}[[leftmargin=*]]%
\item \href{https://www.math.purdue.edu/~yipn/421/2024-11-05-NetFlows.pdf#page=2}{\textbf{Network Flow Problem Formulation}}: $\text{min} \sum_{(i,j) \in \mathcal{A}} C_{ij} X_{ij}$ subject to $\sum_{i} x_{ik} - \sum_{j} x_{kj} = -b_k$ for all $k \in N$ and $X_{ij} \ge 0$ for all $(i,j) \in \mathcal{A}$.<L[2024-11-05-NetFlows] 2><L[2024-11-05-NetFlows] 3>%
\item \href{https://www.math.purdue.edu/~yipn/421/2024-11-05-NetFlows.pdf#page=22}{\textbf{Primal Simplex Method for Network Flows}}: Iteratively identify the entering variable (arc with most negative reduced cost), leaving variable (arc in cycle formed by entering variable and spanning tree), update spanning tree, and update primal/dual variables until all dual slacks are non-negative.<L[2024-11-05-NetFlows] 22>%
\item \href{https://www.math.purdue.edu/~yipn/421/2024-11-05-NetFlows.pdf#page=15}{\textbf{Finding a Spanning Tree Basis}}: A square submatrix of $\tilde{A}$ (incidence matrix without root node row) is a basis if and only if the arcs corresponding to its columns form a spanning tree.<L[2024-11-05-NetFlows] 15><L[2024-11-05-NetFlows] 16>%
\item \href{https://www.math.purdue.edu/~yipn/421/2024-11-05-NetFlows.pdf#page=8}{\textbf{Dual Problem Formulation}}: $\text{max} \ -\sum_{i \in N} b_i y_i$ subject to $y_j - y_i + z_{ij} = c_{ij}$ for all $(i,j) \in \mathcal{A}$ and $z_{ij} \ge 0$.<L[2024-11-05-NetFlows] 8><L[2024-11-05-NetFlows] 9>%
\item \href{https://www.math.purdue.edu/~yipn/421/2024-11-05-NetFlows.pdf#page=14}{\textbf{Calculating Primal Flows from a Spanning Tree}}: Iteratively use flow balance equations at leaf nodes to determine flow values, working inward from leaves to root.<L[2024-11-05-NetFlows] 14>%
\item \href{https://www.math.purdue.edu/~yipn/421/2024-11-05-NetFlows.pdf#page=19}{\textbf{Calculating Dual Variables from a Spanning Tree}}: Starting from the root node, iteratively use the dual flow equation $y_j - y_i = c_{ij}$ for arcs in the spanning tree to compute node potentials.<L[2024-11-05-NetFlows] 19>%
\end{itemize}%
\begin{itemize}[[leftmargin=*]]%
\item \href{https://www.math.purdue.edu/~yipn/421/2024-11-07-NetSimplex.pdf#page=2}{\textbf{Network Simplex Method (Primal)}}: Find a spanning tree T. Solve $\tilde{B}\tilde{X} = \tilde{b}$ for primal flow X, where $x_{ij} = 0$ for $(i,j) \notin T$. If infeasible ($x_{ij} < 0$ for some $(i,j)$), iteratively choose an entering arc with $z_{ij} < 0$, find a cycle, and update primal and dual flows using the leaving arc selection rule (smallest flow in reverse direction of entering arc).<L[2024-11-07-NetSimplex] 2><L[2024-11-07-NetSimplex] 5><L[2024-11-07-NetSimplex] 6><L[2024-11-07-NetSimplex] 7><L[2024-11-07-NetSimplex] 8><L[2024-11-07-NetSimplex] 9><L[2024-11-07-NetSimplex] 10><L[2024-11-07-NetSimplex] 11><L[2024-11-07-NetSimplex] 12><L[2024-11-07-NetSimplex] 13><L[2024-11-07-NetSimplex] 14><L[2024-11-07-NetSimplex] 15><L[2024-11-07-NetSimplex] 16><L[2024-11-07-NetSimplex] 17><L[2024-11-07-NetSimplex] 18><L[2024-11-07-NetSimplex] 19><L[2024-11-07-NetSimplex] 20><L[2024-11-07-NetSimplex] 21><L[2024-11-07-NetSimplex] 22><L[2024-11-07-NetSimplex] 23><L[2024-11-07-NetSimplex] 24><L[2024-11-07-NetSimplex] 25><L[2024-11-07-NetSimplex] 26><L[2024-11-07-NetSimplex] 27><L[2024-11-07-NetSimplex] 28><L[2024-11-07-NetSimplex] 30><L[2024-11-07-NetSimplex] 31>%
\end{itemize}%
\begin{itemize}[[leftmargin=*]]%
\item \href{https://www.math.purdue.edu/~yipn/421/2024-11-14-NetAppls.pdf#page=1}{\textbf{Transportation Problem}}: $\sum_{j} x_{ij} \le s_i$ and $\sum_{i} x_{ij} \ge d_j$ <L[2024-11-14-NetAppls] 1><L[2024-11-14-NetAppls] 2>%
\item \href{https://www.math.purdue.edu/~yipn/421/2024-11-14-NetAppls.pdf#page=3}{\textbf{Transportation Problem with Dummy Node}}: $\sum_{j} x_{ij} + x_{id} = s_i$ and $\sum_{i} x_{ij} - d_j + x_{d_j} = d_j$ <L[2024-11-14-NetAppls] 3>%
\item \href{https://www.math.purdue.edu/~yipn/421/2024-11-14-NetAppls.pdf#page=4}{\textbf{Balance Condition}}: $\sum_{i} s_i - \sum_{i} x_{id} = \sum_{j} d_j + \sum_{j} x_{d_j}$ <L[2024-11-14-NetAppls] 4>%
\item \href{https://www.math.purdue.edu/~yipn/421/2024-11-14-NetAppls.pdf#page=7}{\textbf{General Inequality Constraints}}: $A_1 \bar{x} \ge -\bar{b}_1$, $A_2 \bar{x} \le -\bar{b}_2$, $A_1 \bar{x} = -\bar{b}_1 + \bar{w}_1$, $A_2 \bar{x} = -\bar{b}_2 - \bar{w}_2$, $\sum_{p} w_{1p} - \sum_{s} w_{s2} = \sum_{p} b_{1p} + \sum_{s} b_{s2}$ <L[2024-11-14-NetAppls] 7>%
\item \href{https://www.math.purdue.edu/~yipn/421/2024-11-14-NetAppls.pdf#page=9}{\textbf{Upper Bounded Transshipment Problem}}: $AX = b$, $0 \le X \le U$ <L[2024-11-14-NetAppls] 9><L[2024-11-14-NetAppls] 10><L[2024-11-14-NetAppls] 11><L[2024-11-14-NetAppls] 12>%
\item \href{https://www.math.purdue.edu/~yipn/421/2024-11-14-NetAppls.pdf#page=17}{\textbf{Bellman's Equation}}: $v_i = \min_{j} \{c_{ij} + v_j\}$, $v_r = 0$ <L[2024-11-14-NetAppls] 17>%
\item \href{https://www.math.purdue.edu/~yipn/421/2024-11-14-NetAppls.pdf#page=18}{\textbf{Bellman{-}Ford Algorithm}}: $k = 0$, $v_r^{(0)} = 0$, $v_i^{(0)} = +\infty$, $i \neq r$; $k \ge 0$, $v_r^{(k)} = 0$, $v_i^{(k+1)} = \min_j \{c_{ij} + v_j^{(k)}\}$, $i \neq r$ <L[2024-11-14-NetAppls] 18><L[2024-11-14-NetAppls] 19>%
\end{itemize}%
\begin{itemize}[[leftmargin=*]]%
\item \href{https://www.math.purdue.edu/~yipn/421/2024-11-19-MaxFlowMinCut.pdf#page=4}{\textbf{Max Flow Min Cut Theorem}}: The maximum flow from source to sink in a network is equal to the minimum capacity of any cut separating the source and sink.<L[2024-11-19-MaxFlowMinCut] 4>%
\item \href{https://www.math.purdue.edu/~yipn/421/2024-11-19-MaxFlowMinCut.pdf#page=10}{\textbf{Ford{-}Fulkerson Algorithm}}: Iteratively finds augmenting paths in a network and updates the flow along these paths until no more augmenting paths exist.<L[2024-11-19-MaxFlowMinCut] 10><L[2024-11-19-MaxFlowMinCut] 11><L[2024-11-19-MaxFlowMinCut] 12><L[2024-11-19-MaxFlowMinCut] 13>%
\item \href{https://www.math.purdue.edu/~yipn/421/2024-11-19-MaxFlowMinCut.pdf#page=14}{\textbf{Linear Programming Formulation of Max Flow}}: Minimize $-x_{ts}$ subject to flow conservation constraints at each node and capacity constraints on each edge, where $x_{ts}$ represents the flow from source to sink through a fictitious arc with infinite capacity.<L[2024-11-19-MaxFlowMinCut] 14><L[2024-11-19-MaxFlowMinCut] 15>%
\end{itemize}%
\begin{itemize}[[leftmargin=*]]%
%
\end{itemize}%
\begin{itemize}[[leftmargin=*]]%
\item \href{https://www.math.purdue.edu/~yipn/421/2024-12-03-Duality.pdf#page=1}{\textbf{Primal Problem}}: $\text{max } c^Tx \text{ s.t. } Ax \le b, x \ge 0$ <L[2024-12-03-Duality] 1>%
\item \href{https://www.math.purdue.edu/~yipn/421/2024-12-03-Duality.pdf#page=1}{\textbf{Dual Problem}}: $\text{min } b^Ty \text{ s.t. } A^Ty \ge c, y \ge 0$ <L[2024-12-03-Duality] 1>%
\item \href{https://www.math.purdue.edu/~yipn/421/2024-12-03-Duality.pdf#page=2}{\textbf{Complementary Slackness}}: $c^Tx = y^TAx = y^Tb$; $0 = (y^TA - c^T)x$; $0 = y^T(Ax - b)$; $x_jz_j = 0$; $y_iw_i = 0$ <L[2024-12-03-Duality] 2>%
\item \href{https://www.math.purdue.edu/~yipn/421/2024-12-03-Duality.pdf#page=3}{\textbf{Graphical Method}}: Solve a two-variable LP problem by plotting the constraints to define a feasible region and then finding the vertex that maximizes the objective function. <L[2024-12-03-Duality] 3><L[2024-12-03-Duality] 4>%
\item \href{https://www.math.purdue.edu/~yipn/421/2024-12-03-Duality.pdf#page=14}{\textbf{Sensitivity Analysis (c changes)}}: $\Delta \mathcal{S} = (\Delta c)^T x^$ <L[2024-12-03-Duality] 14><L[2024-12-03-Duality] 15><L[2024-12-03-Duality] 20><L[2024-12-03-Duality] 21>%
\item \href{https://www.math.purdue.edu/~yipn/421/2024-12-03-Duality.pdf#page=15}{\textbf{Sensitivity Analysis (b changes)}}: $\Delta \mathcal{S} = (\Delta b)^T y^$ <L[2024-12-03-Duality] 15><L[2024-12-03-Duality] 20><L[2024-12-03-Duality] 21>%
\item \href{https://www.math.purdue.edu/~yipn/421/2024-12-03-Duality.pdf#page=15}{\textbf{Shadow Prices}}: $y^$ represents the marginal value of a resource. <L[2024-12-03-Duality] 15><L[2024-12-03-Duality] 20><L[2024-12-03-Duality] 21>%
\end{itemize}%
\newpage

%
\end{document}